% Created 2026-09-18 Fri 20:28
% Intended LaTeX compiler: xelatex
\documentclass[11pt, a4paper, abstract=yes]{scrartcl}
\usepackage{hyperref}
\usepackage{amsthm,amssymb,amsmath}
\usepackage{hyperref}
\usepackage{graphicx}
\usepackage{iftex}
\ifPDFTeX
\usepackage[T1]{fontenc}\usepackage{lmodern}
\else
\usepackage{fontspec}
\IfFontExistsTF{Latin Modern Roman}{\setmainfont{Latin Modern Roman}}{}
\IfFontExistsTF{Latin Modern Sans}{\setsansfont{Latin Modern Sans}}{}
\IfFontExistsTF{Latin Modern Mono}{\setmonofont{Latin Modern Mono}}{}
\fi
\theoremstyle{definition}
\newtheorem{definition}{Definition}
\newtheorem{mexample}[definition]{Example}
\theoremstyle{remark}
\newtheorem{remark}[definition]{Remark}
\theoremstyle{plain}
\newtheorem{lemma}[definition]{Lemma}
\newtheorem{claim}[definition]{Claim}
\newtheorem{proposition}[definition]{Proposition}
\newtheorem{theorem}[definition]{Theorem}
\newtheorem{corollary}[definition]{Corollary}
\newtheorem{cor}[definition]{Corollary}
\newtheorem{conjecture}[definition]{Conjecture}
\newtheorem{question}[definition]{Question}
\newcommand{\set}[1]{\left\{{#1}\right\}}
\newcommand\setsuchas[2]{\left\{\,{#1}\,\vrule\,{#2}\,\right\}}
\newcommand{\vek}[1]{\boldsymbol{#1}}
\newcommand{\NN}{{\mathbb{N}}}
\newcommand{\NNp}{{\mathbb{N}_+}}
\newcommand{\ZZ}{{\mathbb{Z}}}
\newcommand{\CC}{{\mathbb{C}}}
\newcommand{\RR}{{\mathbb{R}}}
\newcommand{\QQ}{{\mathbb{Q}}}
\newcommand{\seqnum}[1]{\href{https://oeis.org/#1}{\textrm \underline{#1}}}
\author{Jan Snellman\thanks{Department of Mathematics, Linköping University, 581 83 Linköping, Sweden. \texttt{jan.snellman@liu.se}. ORCID: \href{https://orcid.org/0009-0002-6676-5068}{0009-0002-6676-5068}}}
\usepackage[notcite]{showkeys}
\usepackage{iftex}
\ifPDFTeX\usepackage[T1]{fontenc}\else\usepackage{fontspec}\fi
\usepackage{orcidlink}
\usepackage{booktabs}
\usepackage{etoolbox}
\AtEndPreamble{\AtEveryBibitem{\clearfield{note}\clearfield{annotation}\clearfield{addendum}}}
\hypersetup{hidelinks}
\newcommand{\Koz}{\operatorname{Koz}}
\newcommand{\shift}{\operatorname{shift}}
\newcommand{\ff}{\operatorname{ff}}
\newcommand{\inn}{\operatorname{in}}
\newcommand{\conv}{\operatorname{conv}}
\newcommand{\indeg}{\operatorname{indeg}}
\newcommand{\tp}{\operatorname{top}}
\newcommand{\HF}{\operatorname{HF}}
\author{Jan Snellman \orcidlink{0009-0002-6676-5068}}
\date{\today}
\title{Hilbert functions over the exterior algebra and \(f\)-vectors in higher rank\\\medskip
\large II: Generalizing Linusson's count to sumsets of f-vectors}
\hypersetup{
 pdfauthor={Jan Snellman \orcidlink{0009-0002-6676-5068}},
 pdftitle={Hilbert functions over the exterior algebra and \(f\)-vectors in higher rank},
 pdfkeywords={},
 pdfsubject={},
 pdfcreator={Emacs 30.2 (Org mode 9.7.11)}, 
 pdflang={English}}
\usepackage[doi=false,url=false,isbn=false,style=alphabetic]{biblatex}
\addbibresource{../literature/exterior-convex-references.bib}
\begin{document}

\maketitle
\begin{abstract}
Let \(E=\bigwedge(k^n)\) be the exterior algebra on \(n\) generators, let
\(F=\bigoplus_{i=1}^{r}Eg_i\) be a graded free \(E\)-module with \(\deg g_i=d_i\), and let \(M\)
range over the monomial submodules of Amata and Crupi. The Hilbert function of \(F/M\) is a shifted
sum of \(f\)-vectors of simplicial complexes, and the companion paper \autocite{SnellmanArticle1} determines the convex hull of the
set of such sums. Here we count them. For \(r=1\) the count is Linusson's (1999), obtained from a
recursion through compressed complexes; we show that \textbf{no new recursion is needed for general
\(r\)}. Introducing two statistics of an \(f\)-vector --- its top nonvanishing degree and the initial
degree of the corresponding ideal --- we prove that the number of achievable Hilbert functions is a
product of \(r-1\) transfer matrices, one per consecutive generator-degree gap, whose entries are
differences of Linusson's own refined counts \(E^p(n,k)\). The rank therefore enters only as the
length of a matrix product. The bridge is a reformulation of the counting problem, proved from
Amata and Crupi's lex submodules, that replaces a sumset cardinality by a count of tuples subject to
a nearest-neighbour constraint. We give the same theorem for the \emph{proper} variant that the companion
paper's polytope is built from, prove the lemma that makes the restriction legitimate ---
lexification preserves properness --- and derive consequences: a linear recurrence in \(r\) for
constant gaps, polynomiality in \(r\) when all gaps vanish, an exact threshold beyond which the count
degenerates to a plain power, and polynomiality in the number of vertices.
\end{abstract}

\medskip\noindent\textbf{Keywords.} exterior algebra; Hilbert function; \(f\)-vector; simplicial
complex; enumeration; transfer matrix; lex submodule; Kruskal--Katona theorem.

\smallskip\noindent\textbf{2020 Mathematics Subject Classification.} Primary 05A15, 13F55;
Secondary 05E40, 13D40, 06A07.
\medskip

\setcounter{tocdepth}{2}
\tableofcontents
\section{Introduction}
\label{sec-intro}
\subsection{The rank-one count}
\label{sec:orgdd6a8db}

Let \(k\) be a field and \(E=E_n=\bigwedge(k^n)\) the exterior algebra. Every monomial of \(E\) is
squarefree, so every monomial ideal of \(E\) is the Stanley--Reisner ideal \autocite{Stanley1996} of an order
family \(\Delta\subseteq2^{[n]}\), and \(\dim_k(E/I_\Delta)_j=f_j(\Delta)\) exactly. Since passing to
an initial ideal preserves Hilbert functions, the set of Hilbert functions of graded quotients of
\(E\) \emph{is} the set of \(f\)-vectors of simplicial complexes on at most \(n\) vertices.

Three classical results describe that set. Kruskal--Katona \autocite{Kruskal1963,Katona1968} determines it exactly; Kozlov \textcite{Kozlov1997}
computes its convex hull; and Linusson (1999) \textbf{counts} it. Linusson's method is a recursion on
compressed complexes: peel off the last vertex, split the complex into the part avoiding it and the
link of it, and record how far the first part is filled. The bookkeeping variable of that recursion
--- his \(E^p(n,k)\), the number of \(f\)-vectors that are maximal up to level \(k\) and no further
--- is the object this paper turns out to need.
\subsection{The rank-\(r\) question}
\label{sec:org959e367}

Let \(F=\bigoplus_{i=1}^rEg_i\) with \(\deg g_i=d_i\in\mathbb Z\), \(d_1\le\cdots\le d_r\). Amata and
Crupi's monomial submodules are those generated by monomials of \(F\); such an \(M\) is automatically
a direct sum \(I_1g_1\oplus\cdots\oplus I_rg_r\) of monomial ideals of \(E\), and the Hilbert function
of the quotient splits,
\begin{equation}\label{eq:auto-1}
  \HF(F/M)(j)\;=\;\sum_{i=1}^{r}f_{j-d_i}(\Delta_i),
\end{equation}
a sum of \(r\) independently chosen \(f\)-vectors, each shifted by its own generator degree. Writing
\(\mathcal F_n\) for the set of \(f\)-vectors of order ideals of \(\mathcal P([n])\) and
\(\shift^d\) for the degree-shift, the achievable set is the sumset
\begin{equation}\label{eq:auto-2}
  \mathrm{FF}(\vec d)\;=\;\Bigl\lbrace\sum_{i=1}^{r}\shift^{d_i}(\vec s_i)\;:\;\vec s_i\in\mathcal F_n\Bigr\rbrace .
\end{equation}
The companion paper \textcite{SnellmanArticle1} determines \(\conv(\mathrm{FF}(\vec d))\): it is a Minkowski sum of \(r\) shifted copies
of the Kozlov simplex. \textbf{This paper computes \(|\mathrm{FF}(\vec d)|\).}

The problem is a sumset cardinality, and those are usually hard. Two facts show where the difficulty
sits. The map \((\vec s_1,\dots,\vec s_r)\mapsto\sum_i\shift^{d_i}(\vec s_i)\) is injective exactly when the
generator degrees are far enough apart --- precisely, when every gap \(d_{i+1}-d_i\) is at least
\(n+1\), and in that regime \(|\mathrm{FF}(\vec d)|=|\mathcal F_n|^r\) and there is nothing to do.
Everything of interest happens when the shifts overlap and distinct tuples collide.
\subsection{The result}
\label{sec:orgc502b9c}

The answer is a transfer-matrix product, and it is built from Linusson's own refinement rather than
from a new recursion. Attach to \(\vec s\in\mathcal F_n\) two statistics: \(\tp(\vec s)\), the greatest \(j\)
with \(s(j)>0\); and \(\indeg(\vec s)\), the least \(j\) with \(s(j)<\binom nj\), the initial degree
of the corresponding monomial ideal. Let \(A[\tau][\rho]\) count the \(\vec s\in\mathcal F_n\) with
\(\tp(\vec s)=\tau\) and \(\indeg(\vec s)=\rho\). Then (Theorem \ref{thm-main})
\begin{equation}\label{eq:auto-3}
  \lvert\mathrm{FF}(\vec d)\rvert\;=\;\vec u\cdot T(\delta_2)\,T(\delta_3)\cdots T(\delta_r)\cdot\vec 1,
  \qquad \delta_i=d_i-d_{i-1},
\end{equation}
where \(T(\delta)[\tau][\tau']=\sum_{\rho>\tau-\delta}A[\tau'][\rho]\) and \(\vec u\) is the \(\tp\)
marginal of \(A\). And \(A\) is not a new object: Lemma \ref{lem-joint} identifies it as a difference of
Linusson's numbers,
\begin{equation}\label{eq:auto-4}
  A[\tau][\rho]\;=\;E^n(\tau,\rho-1)-E^n(\tau-1,\rho-1),
\end{equation}
so one run of his recursion at \(p=n\) supplies everything. The rank \(r\) enters only as the length
of a matrix product.

The bridge is Theorem \ref{thm-chaining}, which replaces the sumset by a count of tuples subject to a
constraint relating consecutive entries only:
\begin{equation}\label{eq:auto-5}
  \lvert\mathrm{FF}(\vec d)\rvert=\#\bigl\lbrace(\vec s_1,\dots,\vec s_r)\in\mathcal F_n^{\,r}:
   \tp(\vec s_{i-1})<\indeg(\vec s_i)+\delta_i\ \text{ for } i=2,\dots,r\bigr\rbrace .
\end{equation}
That is proved from Amata and Crupi's theory of lex submodules \autocite{AmataCrupi2020,AmataCrupi2020KK}, which supplies a canonical
representative for each achievable Hilbert function; the constraint is their Definition 2.7 chaining
condition read off the \(f\)-vectors. Once the constraint is nearest-neighbour, the transfer-matrix
form is a rearrangement.
\subsection{What is proved, and what is only observed}
\label{sec:orgfc4a2f1}

The distinction is kept explicit throughout, and the reader is asked to take the labels seriously.

\begin{center}
\small
\begin{tabular}{ll}
result & status\\
\hline
Thm. \ref{thm-linusson}, Linusson's recursion & cited; transcription verified against his tables\\
Lem. \ref{lem-joint}, the joint distribution & proved; verified entrywise for \(n\le6\)\\
Thm. \ref{thm-chaining}, the chaining reformulation & proved\\
Thm. \ref{thm-main}, the counting theorem & proved; verified, 41 configurations\\
Lem. \ref{lem-lexproper}, lexification preserves properness & proved\\
Thm. \ref{thm-proper}, the proper count & proved; reproduces recorded sequences\\
Prop. \ref{prop-recurrence}, recurrence in \(r\) & proved\\
Prop. \ref{prop-delta0}, degree bound at \(\delta=0\) & proved\\
Thm. \ref{thm-delta0-degrees}, the \textbf{exact} degrees at \(\delta=0\) & proved; verified, \(n\le8\)\\
Prop. \ref{prop-perron-threshold}, the Perron root above threshold & proved; verified, \(n\le6\)\\
Prop. \ref{prop-normalform}, the normal form \(T(\delta)=P^{\delta}B\) & proved; verified, \(n\le8\)\\
Thm. \ref{thm-poly-p}, polynomiality in \(p\) & proved modulo Linusson's Thm. 5.3\\
Thm. \ref{thm-asymptotics-p}, degree and leading coefficient in \(p\) & proved modulo the same; verified\\
Thm. \ref{thm-asymptotics-p-proper}, the proper analogue in \(p\) & proved modulo the same; verified\\
Prop. \ref{prop-fnn2}, \(\lvert S_n\rvert=F^n(n-2)\) & proved, by Alexander duality\\
Cor. \ref{cor-persymmetry}, \(A\) is persymmetric & proved\\
\end{tabular}
\end{center}

The asymptotics in \(n\), the analogue of Linusson's own \S2.2, are not attempted.
\section{Notation}
\label{sec-notation}
Collected here, before anything uses it. Pointers are to the defining statement; nothing in this
section is needed to read the introduction. The companion article carries a table in the same spirit,
and where the two papers share a symbol they share its meaning.

\begingroup\small
\begin{center}
\begin{tabular}{@{}ll@{}}
  \toprule
  \multicolumn{2}{@{}l}{\emph{Complexes and their \(f\)-vectors}}\\
  \(\mathcal P([n])\)              & the lattice of subsets of \([n]\) \\
  \(\Delta\)                       & an order ideal of \(\mathcal P([n])\); the void one is allowed \\
  \(f_j(\Delta)\)                  & the number of \(j\)-element members of \(\Delta\) \\
  \(\mathcal F_n\)                 & the \(f\)-vectors of all order ideals, Definition~\ref{def-fn} \\
  \(S_n\)                          & the \emph{proper} ones, \(f_0=1\) and \(f_1=n\) \\
  \(\shift^d\)                     & prepends \(d\) zeros \\
  \addlinespace
  \multicolumn{2}{@{}l}{\emph{The two statistics}}\\
  \(\tp(\vec s)\)                  & the largest \(j\) with \(s(j)\neq0\), Definition~\ref{def-stats} \\
  \(\indeg(\vec s)\)               & the least \(j\) with \(s(j)\) below the binomial bound \\
  \(\indeg(I)\)                    & of a monomial ideal, Definition~\ref{def-setup} \\
  \(A[\tau][\rho]\)                & the joint distribution of \((\tp,\indeg)\), Lemma~\ref{lem-joint} \\
  \addlinespace
  \multicolumn{2}{@{}l}{\emph{Sumsets and the count}}\\
  \(\vec d=(d_1,\dots,d_r)\)       & the shift vector, sorted; \(\delta_i=d_i-d_{i-1}\) its gaps \\
  \(\mathrm{FF}(\vec d)\), \(\mathrm{FF}_{\mathrm{pr}}(\vec d)\) & the sumsets, Definition~\ref{def-fn} \\
  \(F^p(n)\), \(E^p(n,k)\)         & Linusson's counts, Definition~\ref{def-linusson} \\
  \(\vec u\), \(T(\delta)\), \(\vec 1\) & row vector, transfer matrix, column of ones, Theorem~\ref{thm-main} \\
  \bottomrule
\end{tabular}
\end{center}
\endgroup

Two points of contact with the companion article are worth making explicit, since the same objects
are handled there under a different convention.

First, \textbf{\textbf{that article distinguishes three conventions}} for "simplicial complex on the ground set
\([n]\)": (i) all order ideals of \(\mathcal P([n])\), (ii) the non-empty ones, (iii) those containing
\(\emptyset\) and every singleton. \textbf{This paper works throughout in convention (i)} --- the void family
is a member of \(\mathcal F_n\), and it has to be, since \(\indeg\) attains its minimum \(0\) there and
that is what fixes the threshold of Proposition \ref{prop-threshold}. Convention (iii) is what \(S_n\)
picks out, and the two papers agree on the symbol \(S_n\).

Second, the companion article writes \(S_n^{+}\) for the \(f\)-vectors of the \textbf{non-empty} order
ideals, convention (ii). In the notation here that is
\begin{equation}\label{eq:Fn-vs-Splus}
  \mathcal F_n \;=\; S_n^{+}\cup\lbrace\vec 0\rbrace,\qquad
  \lvert\mathcal F_n\rvert \;=\; \lvert S_n^{+}\rvert+1 .
\end{equation}
We keep \(\mathcal F_n\) rather than \(S_n^{+}\cup\lbrace\vec0\rbrace\) because every matrix below is
indexed by it.
\section{Setup}
\label{sec-setup}
We recall only what is needed; the companion paper has the full development.

\begin{definition}
\(E=E_n=\bigwedge(k^n)=k\langle e_1,\dots,e_n\rangle/(e_ie_j+e_je_i,\ e_i^2)\), with
\(\mathfrak m=(e_1,\dots,e_n)\), and for \(\sigma\subseteq[n]\), \(e_\sigma\) the corresponding
monomial. For \(r\ge1\) and \(d_1\le\cdots\le d_r\) in \(\mathbb Z\), let
\(F=\bigoplus_{i=1}^rEg_i\) with \(\deg g_i=d_i\). A \textbf{monomial submodule} is one generated by
monomials \(e_\sigma g_i\); such an \(M\) decomposes uniquely as \(M=\bigoplus_iI_ig_i\) with each
\(I_i\subseteq E\) a monomial ideal \textcite[Definition 2.1]{AmataCrupi2020}. It is \textbf{proper} if no \(I_i\) has
a minimal generator of degree \(\le1\).

For a monomial ideal \(I\subseteq E\) write
\begin{equation}\label{eq:indeg-ideal}
  \indeg(I)=\min\lbrace j:I_j\ne0\rbrace ,
\end{equation}
its \textbf{initial degree}, with \(\indeg(0)=n+1\) since the minimum is then over an empty set.
Equivalently it is the least degree of a minimal generator of \(I\): if \(j_0\) is least with
\(I_{j_0}\ne0\) then \((\mathfrak mI)_{j_0}\) is built from \(I_{j_0-1}=0\), so every non-zero
element of \(I_{j_0}\) is a minimal generator, and no minimal generator has smaller degree.
\label{def-setup}
\end{definition}

\begin{definition}
An \textbf{order ideal} of \(\mathcal P([n])\), the lattice of subsets of \([n]\), is a family
\(\Delta\subseteq\mathcal P([n])\) closed under taking subsets, with
\(f_j(\Delta)=|\lbrace\sigma\in\Delta:|\sigma|=j\rbrace|\). Put
\begin{equation}\label{eq:auto-6}
  \mathcal F_n=\lbrace(f_0(\Delta),\dots,f_n(\Delta)):\Delta\ \text{an order ideal of }\mathcal P([n])\rbrace,
\end{equation}
the void family included and singletons not required, and let
\(S_n=\lbrace \vec s\in\mathcal F_n:s(0)=1,\ s(1)=n\rbrace\) be the \textbf{proper} elements --- the \(f\)-vectors
of simplicial complexes using the full vertex set. For \(d\ge0\), \(\shift^d\) prepends \(d\) zeros.
Finally
\begin{equation}\label{eq:auto-7}
  \mathrm{FF}(\vec d)=\Bigl\lbrace\sum_i\shift^{d_i}(\vec s_i):\vec s_i\in\mathcal F_n\Bigr\rbrace,\qquad
  \mathrm{FF}_{\mathrm{pr}}(\vec d)=\Bigl\lbrace\sum_i\shift^{d_i}(\vec s_i):\vec s_i\in S_n\Bigr\rbrace .
\end{equation}
\label{def-fn}
\end{definition}
\subsection{Notation: which quantities are tuples}
\label{sec:orgd2532ee}

\textbf{An arrow marks a variable whose value is a tuple; a plain letter is a scalar.} The same convention
is used in the companion article, and it is worth stating in full because two different kinds of
index occur here and the difference matters.

\begin{itemize}
\item \(\vec s\in\mathcal F_n\) is a whole \(f\)-vector; \(\vec d=(d_1,\dots,d_r)\) is the shift vector;
\(\vec u\) and \(\vec 1\) are the row and column vectors of Theorem \ref{thm-main}. Each is a tuple, and
each carries an arrow.
\item \textbf{A plain subscript on an arrowed letter selects a member of the \(r\)-tuple}, and that member may
itself be a vector or a scalar: \(\vec s_i\) is the \(i\)-th \(f\)-vector of an \(r\)-tuple, and so
still carries an arrow; \(d_i\) is the \(i\)-th generator degree, a scalar, and does not.
\item \textbf{A parenthesised argument selects a coordinate}, and is always a scalar: \(s(j)\) is the number of
\(j\)-subsets counted by \(\vec s\), \(\HF(F/M)(j)\) is a dimension, \(x(j)\) a single entry of
\(\vec x\). So \(s_i(j)\) (arrowed letter, tuple subscript, coordinate argument) is the
\(j\)-th coordinate of the \(i\)-th \(f\)-vector, and is a scalar.
\item Matrices are capitals and never carry an arrow (\(A\), \(T(\delta)\)); their entries are written
\(A[\tau][\rho]\), \(u[\tau]\), also scalars.
\item \textbf{Exception, and it is the usual one:} a subscript on a module or a graded object still means its
homogeneous component: \(E_j\), \(F_j\), \(I_j\), \(\mathrm{Mon}_j(L)\) are subspaces, not
coordinates of a vector, and follow the standard convention rather than this one.
\end{itemize}

The one place the distinction is load-bearing is Definition \ref{def-fn} against
Theorem \ref{thm-chaining}: without it, \(s_1\) would mean both "the first coordinate of \(\vec s\)"
(where it equals \(n\) for a proper vector) and "the first \(f\)-vector of the tuple". Under the
convention those are \(s(1)\) and \(\vec s_1\), and cannot be confused.

The dictionary between the algebra and the combinatorics is exact: \(I\mapsto\Delta(I)=\lbrace\sigma:
e_\sigma\notin I\rbrace\) is a bijection from monomial ideals of \(E\) onto order ideals of \(\mathcal P([n])\),
\(\HF(E/I)\) is the \(f\)-vector of \(\Delta(I)\), and both degenerate cases are included
(\(I=E\mapsto\) the void family, \(I=0\mapsto\) the full simplex). Consequently
\(\mathrm{FF}(\vec d)=\lbrace\HF(F/M):M\ \text{monomial}\rbrace\) and
\(\mathrm{FF}_{\mathrm{pr}}(\vec d)=\lbrace\HF(F/M):M\ \text{proper monomial}\rbrace\).

\begin{remark}
Both sets depend on \(\vec d\) only through the gaps \(\delta_i=d_i-d_{i-1}\ge0\): a global shift
re-indexes \(\mathrm{FF}(\vec d)\) bijectively. Negative generator degrees are therefore allowed
throughout, as they are in Amata and Crupi. What is \textbf{not} optional is that \(\vec d\) be sorted; see
Remark \ref{rem-sorted}.
\label{rem-shift-only-gaps}
\end{remark}
\section{Linusson's recursion}
\label{sec-linusson}
\begin{definition}
\textcite[\S5]{Linusson1999} For \(p\ge0\) and \(n\ge-1\), let \(F^p(n)\) be the number of \(f\)-vectors of
simplicial complexes with \(f_1\le p\) and \(f_j=0\) for \(j>n\), both degenerate \(f\)-vectors
included --- and let \(E^p(n,k)\) count those among them with \(f_i=\binom pi\) for \(i\le k\) but
\(f_{k+1}<\binom p{k+1}\).
\label{def-linusson}
\end{definition}

\begin{theorem}
\textcite[Theorem 5.2]{Linusson1999} For \(0\le k\) and \(1\le n\le p\),
\begin{equation}\label{eq:auto-8}
  F^p(n)=\begin{cases}
    1+\sum_{i=0}^{n}E^{p-1}(n,i)\,F^{p-1}(i-1), & n<p,\\
    1+F^p(p-1), & n=p,
  \end{cases}
\end{equation}
\begin{equation}\label{eq:auto-9}
  E^p(n,k)=\begin{cases}
    \sum_{i=k}^{n}E^{p-1}(n,i)\,E^{p-1}(i-1,k-1), & k\le n<p,\\
    E^p(p-1,k), & k<n=p,\\
    1, & k=n=p,
  \end{cases}
\end{equation}
with \(F^p(-1)=1\), \(F^p(0)=2\), \(F^0(n)=2\), \(E^p(n,-1)=1\), \(E^p(n,n)=1\), and
\(E^p(n,k)=0\) for \(k>n\). Moreover \(F^p(n)=\sum_kE^p(n,k)\).
\label{thm-linusson}
\end{theorem}

\begin{remark}
The transcription above was checked entry by entry against Linusson's own appendix: all \(100\)
printed entries of his Table 5 and all \(99\) of his Table 6 are reproduced exactly, and an
independent brute-force enumeration of order ideals agrees for every \(F^p(n)\) and
\(E^p(n,k)\) with \(p\le5\). One wrinkle is worth recording. At \(k=n=p\) the clause
"\(f_{k+1}<\binom p{k+1}\)" read literally is vacuously false, since \(\binom p{p+1}=0\); but
Linusson's stated boundary value \(E^p(n,n)=1\) and his Table 6 diagonal both require the clause to
be dropped there. We follow Linusson.
\label{rem-transcription}
\end{remark}

Two consequences we use repeatedly. Setting \(p=n\) gives \(F^n(n)=|\mathcal F_n|\); and the
proper count is a difference of two such.

\begin{remark}
\(|S_n|=F^n(n)-F^{n-1}(n-1)\), giving
\begin{equation}\label{eq:auto-10}
  \lvert S_n\rvert = 1,\,2,\,5,\,16,\,70,\,457,\,4908,\,95248,\,3617645,\,286007155
  \qquad(n=1,\dots,10),
\end{equation}
whose first five terms agree with direct enumeration. The sequence \(F^p(p)\) is
OEIS A007695 \autocite{OEIS}. There is also the identity \(|S_n|=F^n(n-2)\), for
instance \(|S_9|=F^9(7)=3617645\), a value printed in Linusson's Table 5 --- which is
Proposition \ref{prop-fnn2} below, proved by Alexander duality. Equivalently, via
\(F^n(n)=1+F^n(n-1)\), it says \(F^n(n-1)-F^n(n-2)=F^{n-1}(n-1)-1\). We do not use it.
\label{obs-fnn2}
\end{remark}
\section{Two statistics, and the joint distribution}
\label{sec-statistics}
\begin{definition}
For \(\vec s\in\mathcal F_n\) put
\begin{equation}\label{eq:auto-11}
  \indeg(\vec s)=\min\lbrace j:s(j)<\tbinom nj\rbrace,\qquad
  \tp(\vec s)=\max\lbrace j:s(j)>0\rbrace,
\end{equation}
with \(\indeg(\vec s)=n+1\) for the full simplex and \(\tp(\vec s)=-1\) for the void family. The
symbol \(\indeg\) is the one of \eqref{eq:indeg-ideal}: this is its combinatorial shadow, and
Lemma \ref{lem-dictionary}(2) is the statement that the two agree.
\label{def-stats}
\end{definition}

\begin{definition}
(\emph{The joint-distribution matrix.}) Let \(A\) be the matrix whose entry
\begin{equation}\label{eq:def-A}
  A[\tau][\rho]=\#\bigl\lbrace \vec s\in\mathcal F_n:\tp(\vec s)=\tau,\ \indeg(\vec s)=\rho\bigr\rbrace
\end{equation}
counts the elements of \(\mathcal F_n\) with the prescribed pair of statistics. \textbf{Rows are indexed by
\(\tp\) and columns by \(\indeg\)}, over
\begin{equation}\label{eq:def-A-range}
  \tau=-1,0,1,\dots,n,\qquad \rho=0,1,\dots,n+1 ,
\end{equation}
so \(A\) is \((n+2)\times(n+2)\). Two features of the range are forced by the conventions of
Definition \ref{def-stats} and are used repeatedly below: \(A\) is supported on \(\rho\le\tau+1\), and
the two extreme rows are the single entries \(A[-1][0]=1\) (the void family) and \(A[n][n+1]=1\)
(the full simplex). Every row sums to the number of \(\vec s\) with that value of \(\tp\), and the
grand total is \(\lvert\mathcal F_n\rvert\).
\label{def-A}
\end{definition}

Note \(\indeg(\vec s)\le\tp(\vec s)+1\) always. For \(\tp(\vec s)<n\) this is because
\(s(\tp(\vec s)+1)=0<\binom n{\tp(\vec s)+1}\), so the minimum defining \(\indeg\) is attained no later
than \(\tp(\vec s)+1\); and for \(\tp(\vec s)=n\) it holds because \(\indeg(\vec s)\le n+1\) outright, by
the convention of Definition \ref{def-stats}. The bound is therefore unconditional, the binomial
coefficient \(\binom n{\tp(\vec s)+1}\) vanishing in exactly the case the second clause covers. So
\(A\) is supported on \(\rho\le\tau+1\).

\begin{lemma}
With the convention \(E^n(-2,\cdot)=0\),
\begin{equation}\label{eq:auto-12}
  A[\tau][\rho]\;=\;E^n(\tau,\rho-1)-E^n(\tau-1,\rho-1).
\end{equation}
\label{lem-joint}
\end{lemma}

\begin{proof}
Work at \(p=n\). Every \(\vec s\in\mathcal F_n\) has \(f_1\le n\); conversely every \(f\)-vector counted by
\(E^n(m,k)\) has \(f_1\le n\), hence belongs to a complex on at most \(n\) vertices, which embeds in
\([n]\). So the objects counted by \(E^n(m,k)\) are exactly certain elements of \(\mathcal F_n\).
Under that identification Linusson's clause "\(f_j=0\) for \(j>m\)" says precisely \(\tp(\vec s)\le m\),
and his clause "\(f_i=\binom ni\) for \(i\le k\) but \(f_{k+1}<\binom n{k+1}\)" (with the boundary
convention of Remark \ref{rem-transcription}) says precisely \(\indeg(\vec s)=k+1\). Hence
\[
  E^n(m,k)=\#\lbrace \vec s\in\mathcal F_n:\tp(\vec s)\le m,\ \indeg(\vec s)=k+1\rbrace ,
\]
and subtracting the count at \(m=\tau-1\) from the count at \(m=\tau\) leaves \(\tp(\vec s)=\tau\)
exactly. For \(\tau=-1\) the subtrahend would count \(f\)-vectors with \(\tp\le-2\), of which there
are none (even the void \(f\)-vector has \(\tp=-1\)), which is the stated convention.
\end{proof}

So one run of Theorem \ref{thm-linusson} at \(p=n\) delivers the entire joint distribution. Its
\(\indeg\) marginal telescopes to \(E^n(n,\rho-1)\), which is Linusson's Table 6.

\begin{mexample}
At \(n=3\) the set \(\mathcal F_3\) has ten elements. Listing them with their statistics --- the
\(f\)-vector, then \((\tp,\indeg)\), gives
\[
\begin{array}{llll}
 (0,0,0,0)\colon(-1,0) & (1,0,0,0)\colon(0,1) & (1,1,0,0)\colon(1,1) & (1,2,0,0)\colon(1,1)\\
 (1,2,1,0)\colon(2,1)  & (1,3,0,0)\colon(1,2) & (1,3,1,0)\colon(2,2) & (1,3,2,0)\colon(2,2)\\
 (1,3,3,0)\colon(2,3)  & (1,3,3,1)\colon(3,4) & &
\end{array}
\]
and collecting them into \eqref{eq:def-A} (rows \(\tau=-1,\dots,3\) downward, columns
\(\rho=0,\dots,4\) rightward)
\[
  A=\begin{pmatrix}
    1&0&0&0&0\\ 0&1&0&0&0\\ 0&2&1&0&0\\ 0&1&2&1&0\\ 0&0&0&0&1
  \end{pmatrix},
  \qquad \text{row sums } (1,1,3,4,1),\ \text{total }10=\lvert\mathcal F_3\rvert .
\]
The support condition \(\rho\le\tau+1\) is the vanishing above the superdiagonal; the two corner
entries are the void family and the full simplex; and the antidiagonal symmetry of
Corollary \ref{cor-persymmetry} is already visible, the \(2\) at \((\tau,\rho)=(1,1)\) matching the
\(2\) at \((2,2)\).
\label{ex-A}
\end{mexample}

\begin{remark}
Lemma \ref{lem-joint} was also checked entrywise, for \(n\le6\) and with no mismatches, against
direct enumeration of the order ideals themselves. For the record, \(A\) at \(n=6\), rows \(\tau=-1,\dots,6\) and
columns \(\rho=0,\dots,7\):
\label{rem-A-verified}
\end{remark}

\begin{equation}\label{eq:auto-13}
\begin{pmatrix}
1&0&0&0&0&0&0&0\\
0&1&0&0&0&0&0&0\\
0&5&1&0&0&0&0&0\\
0&20&14&1&0&0&0&0\\
0&43&84&19&1&0&0&0\\
0&25&143&84&14&1&0&0\\
0&1&25&43&20&5&1&0\\
0&0&0&0&0&0&0&1
\end{pmatrix}
\end{equation}
\subsection{Alexander duality, and the identity of Remark \ref{obs-fnn2}}
\label{sec:org587bd3c}

The matrix just printed is visibly symmetric about its antidiagonal. That is not an accident, and
the same duality proves the identity left open in Section \ref{sec-linusson}.

\begin{lemma}
(\emph{Alexander duality on \(f\)-vectors.}) For an order ideal \(\Delta\) of \(\mathcal P([n])\) put
\(\Delta^{\vee}=\lbrace\sigma:[n]\setminus\sigma\notin\Delta\rbrace\). Then \(\Delta^{\vee}\) is an
order ideal, \((\Delta^{\vee})^{\vee}=\Delta\), and
\begin{equation}\label{eq:alexander-f}
  f_j(\Delta^{\vee})=\tbinom nj-f_{n-j}(\Delta).
\end{equation}
Consequently \(D(\vec s)(j):=\binom nj-s(n-j)\) is an involution of \(\mathcal F_n\), and it conjugates
the two statistics into each other:
\begin{equation}\label{eq:alexander-stats}
  \indeg(D\vec s)=n-\tp(\vec s),\qquad \tp(D\vec s)=n-\indeg(\vec s).
\end{equation}
Of these assertions, only \(\Delta^{\vee}\) being an order ideal --- equivalently, \(D\) mapping
\(\mathcal F_n\) into itself --- uses the combinatorics. Writing
\begin{equation}\label{eq:binomial-box}
  B_n=\bigl\lbrace\vec s\in\mathbb Z_{\ge0}^{n+1}\;:\;s(j)\le\tbinom nj\ \text{ for all }j\bigr\rbrace
  \;\supseteq\;\mathcal F_n ,
\end{equation}
the map \(D\) is an involution of \(B_n\) and satisfies \eqref{eq:alexander-stats} on all of it; the
entrywise bound is all that is needed, and it is needed only for the involution.
\label{lem-alexander}
\end{lemma}

\begin{proof}
If \(\sigma\in\Delta^{\vee}\) and \(\tau\subseteq\sigma\) then \([n]\setminus\sigma\subseteq
[n]\setminus\tau\); were \([n]\setminus\tau\in\Delta\), downward closure would put
\([n]\setminus\sigma\in\Delta\), which it is not. So \(\Delta^{\vee}\) is an order ideal. Double
complementation gives \((\Delta^{\vee})^{\vee}=\Delta\). For \eqref{eq:alexander-f}, complementation
\(\sigma\mapsto[n]\setminus\sigma\) is a bijection from the \(j\)-subsets to the \((n-j)\)-subsets,
under which the \(j\)-members of \(\Delta^{\vee}\) correspond to the \((n-j)\)-subsets \textbf{not} in
\(\Delta\). The formula shows \(D\vec s\) depends only on \(\vec s\), so \(D\) is a well-defined
involution of \(\mathcal F_n\). Finally \(D(\vec s)(j)<\binom nj\) if and only if \(s(n-j)>0\), which
is \eqref{eq:alexander-stats}; the conventions \(\tp=-1\) for the void family and
\(\indeg=n+1\) for the full simplex are exactly what make the two extremes correspond.

For the last paragraph of the statement, observe that this final step used no property of \(\vec s\)
beyond \(0\le s(n-j)\le\binom n{n-j}\): the equivalence \(D(\vec s)(j)<\binom nj\iff s(n-j)>0\) reads
off the definition of \(D\), and dually \(D(\vec s)(j)>0\iff s(n-j)<\binom n{n-j}\). Since
\(\binom nj=\binom n{n-j}\), the same bound gives
\(D(D\vec s)(j)=\binom nj-\bigl(\binom n{n-j}-s(j)\bigr)=s(j)\), so \(D\) is an involution of
\(B_n\). What \eqref{eq:alexander-f} adds is that this involution of the box restricts to one of the
subset \(\mathcal F_n\), and that is where downward closure is used.
\end{proof}

\begin{corollary}
\(A[\tau][\rho]=A[n-\rho][n-\tau]\) for all \(\tau,\rho\): the matrix \(A\) is persymmetric. With
\(A\) read as in \eqref{eq:def-A} for \textbf{arbitrary} integers \(\tau,\rho\) (so that it vanishes
outside the range \eqref{eq:def-A-range}), the identity holds on all of \(\mathbb Z^2\), and no
check that the reflected indices stay in range is needed when it is applied inside a sum.
\label{cor-persymmetry}
\end{corollary}

\begin{proof}
\(D\) is a bijection of \(\mathcal F_n\) carrying the class \((\tp,\indeg)=(\tau,\rho)\) onto the class
\((n-\rho,n-\tau)\), by \eqref{eq:alexander-stats}. For the unrestricted form, note that
\eqref{eq:def-A} counts an empty set whenever \(\tau\) or \(\rho\) leaves the range
\eqref{eq:def-A-range}, and \((\tau,\rho)\mapsto(n-\rho,n-\tau)\) maps that range onto itself, so both
sides vanish together outside it.
\end{proof}

\begin{proposition}
\(\lvert S_n\rvert=F^n(n-2)\) for every \(n\ge1\). The hypothesis is sharp rather than cosmetic: at
\(n=0\) the left side is \(1\), since \(\lbrace\emptyset\rbrace\) is the one order ideal of
\(\mathcal P(\emptyset)\) meeting condition (iii) and its \(f\)-vector is \((1)\), while \(F^0(-2)=0\).
\label{prop-fnn2}
\end{proposition}

\begin{proof}
By Definition \ref{def-linusson} at \(p=n\), where the condition \(f_1\le n\) is automatic,
\(F^n(m)\) is the number of \(\vec s\in\mathcal F_n\) with \(\tp(\vec s)\le m\). Now
\(D(\vec s)(n)=\binom nn-s(0)=1-s(0)\) and \(D(\vec s)(n-1)=\binom n{n-1}-s(1)=n-s(1)\), so
\begin{align*}
  \vec s\in S_n
  &\iff s(0)=1\ \text{and}\ s(1)=n\\
  &\iff D(\vec s)(n)=D(\vec s)(n-1)=0
   \iff \tp(D\vec s)\le n-2 .
\end{align*}
As \(D\) is an involution of \(\mathcal F_n\) by Lemma \ref{lem-alexander}, it restricts to a bijection
from \(S_n\) onto \(\lbrace\vec g\in\mathcal F_n:\tp(\vec g)\le n-2\rbrace\), a set of cardinality
\(F^n(n-2)\).
\end{proof}

In words: \textbf{the proper \(f\)-vectors are Alexander-dual to the \(f\)-vectors of order ideals of
dimension at most \(n-3\)}, and properness, two conditions at the bottom of \(\vec s\), becomes
the vanishing of the top two entries of \(D\vec s\). Including both degenerate \(f\)-vectors in
\(\mathcal F_n\), which Definition \ref{def-fn} does and which convention (i) of Section
\ref{sec-notation} requires, is what makes the duality an involution of the whole set rather than a
partial map: the void family and the full simplex are exchanged, as are \(\lbrace\emptyset\rbrace\)
and the boundary of the simplex.

\begin{remark}
One further consequence, recorded because it constrains any future answer to the bijective question
of Section \ref{sec-verification}. Applying \(D\) to every entry and reversing the order,
\((\vec s_1,\dots,\vec s_r)\mapsto(D\vec s_r,\dots,D\vec s_1)\) carries tuples chained for
\((\delta_2,\dots,\delta_r)\) to tuples chained for the \textbf{reversed} gap vector
\((\delta_r,\dots,\delta_2)\), by \eqref{eq:alexander-stats}. So \(|\mathrm{FF}(\vec d)|\) is
invariant under reversing the gaps --- verified for \(n\le6\) over six gap vectors, including
non-palindromic ones. This is an ingredient for a bijective proof of
Theorem \ref{thm-chaining} and a symmetry any such proof must respect, not a solution to that problem.

Two things delimit the symmetry sharply, both settled by computation, and the second of them is the
subject of Remark \ref{rem-proper-not-reversible} below.

\textbf{Nothing more than reversal is true:} \(\lvert\mathrm{FF}\rvert\) is not invariant under arbitrary
permutation of the gaps. At \(n=2\), \(r=4\), the three arrangements of the multiset
\(\lbrace0,0,1\rbrace\) give
\begin{equation}\label{eq:no-permutation-invariance}
  (0,0,1)\mapsto96,\qquad (0,1,0)\mapsto111,\qquad (1,0,0)\mapsto96 ,
\end{equation}
and the two that agree are precisely the two that are reverses of each other. So the involution of
this remark is the whole of the symmetry, not a visible part of a larger one.
\label{rem-gap-reversal}
\end{remark}

\begin{remark}
(\emph{Gap reversal fails for the proper count.}) The second delimitation, and the reason the proper
variant of Section \ref{sec-proper} is not simply the improper theory with a restriction imposed: it
does not inherit the symmetry of Remark \ref{rem-gap-reversal}.

The argument there cannot be run for \(\mathrm{FF}_{\mathrm{pr}}(\vec d)\), because \(D\) does not
preserve \(S_n\): by Proposition \ref{prop-fnn2} it carries \(S_n\) onto
\(\lbrace\vec g\in\mathcal F_n:\tp(\vec g)\le n-2\rbrace\), a different subset of \(\mathcal F_n\), so
the reversal map of Remark \ref{rem-gap-reversal} does not send proper-chained tuples to proper-chained
tuples.

That much would leave open whether the invariance nonetheless holds for some other reason. It does
not: the conclusion itself is false, not merely unproved. At \(n=3\) the gap vectors
\((\delta_2,\delta_3)=(0,1)\) and \((1,0)\) give
\begin{equation}\label{eq:proper-not-reversible}
  \lvert\mathrm{FF}_{\mathrm{pr}}(3;0,0,1)\rvert=45,
  \qquad
  \lvert\mathrm{FF}_{\mathrm{pr}}(3;0,1,1)\rvert=51,
\end{equation}
while the corresponding improper counts agree at \(236\), exactly as Remark
\ref{rem-gap-reversal} predicts. So \(45\ne51\) is a counterexample and not a gap in an argument, and
any attempt to transport the reversal symmetry to the proper count must fail rather than merely
require more work.

The asymmetry is worth keeping in view when reading Section \ref{sec-proper}. Properness is two
conditions at the bottom of \(\vec s\), and Alexander duality moves the bottom of a vector to its
top; a symmetry built out of that duality therefore has no reason to respect a condition pinned to
one end, and here it does not.
\label{rem-proper-not-reversible}
\end{remark}
\section{The chaining reformulation}
\label{sec-chaining}
This section replaces a sumset cardinality by a count of tuples with a nearest-neighbour constraint.
The mechanism is that Amata and Crupi's lex submodules give a canonical representative for each
achievable Hilbert function, and that being a lex submodule is a condition on consecutive components
only.

\begin{lemma}
Let \(I\subseteq E\) be a monomial ideal and \(\vec s=\HF(E/I)\in\mathcal F_n\).
\begin{enumerate}
\item For every integer \(t\ge0\): \(\mathfrak m^t\subseteq I\) if and only if \(\tp(\vec s)<t\).
\item \(\indeg(I)=\indeg(\vec s)\).
\end{enumerate}
\label{lem-dictionary}
\end{lemma}

\begin{proof}
(1) \(\mathfrak m^t\) is spanned by the monomials of degree \(\ge t\), so
\(\mathfrak m^t\subseteq I\) iff \(I_j=E_j\) for all \(j\ge t\), iff \(s(j)=0\) for all \(j\ge t\),
iff \(\tp(\vec s)<t\). This covers \(t=0\) (then \(\mathfrak m^0=E\subseteq I\) iff \(I=E\) iff
\(\tp(\vec s)=-1<0\)) and \(t>n\) (then \(\mathfrak m^t=0\) and \(\tp(\vec s)\le n<t\), both sides true). For
\(t<0\) the equivalence fails, the left side being vacuously true and the right side impossible ---
which is exactly why sortedness of \(\vec d\) is needed in Theorem \ref{thm-chaining}.
(2) \(\dim I_j=\binom nj-s(j)\), so the least \(j\) with \(I_j\ne0\) is the least \(j\) with
\(s(j)<\binom nj\). For \(I=0\) both sides are \(n+1\) by convention; any value exceeding \(n\) gives
the same truth value in (1), so the convention is harmless.
\end{proof}

\begin{theorem}
Fix \(n\ge1\), \(r\ge1\) and integers \(d_1\le d_2\le\cdots\le d_r\), with
\(\delta_i=d_i-d_{i-1}\). Then
\begin{equation}\label{eq:auto-14}
  \lvert\mathrm{FF}(\vec d)\rvert=\#\bigl\lbrace(\vec s_1,\dots,\vec s_r)\in\mathcal F_n^{\,r}\;:\;
  \tp(\vec s_{i-1})<\indeg(\vec s_i)+\delta_i\ \text{ for } i=2,\dots,r\bigr\rbrace .
\end{equation}
\label{thm-chaining}
\end{theorem}

\begin{proof}
(a) By Definition \ref{def-setup} and the dictionary, \(\mathrm{FF}(\vec d)\) is the set of \(\HF(F/M)\)
for \(M\) a monomial submodule. Indeed, for \(M=\bigoplus_iI_ig_i\),
\[
  \HF(F/M)(j)=\sum_i\dim(E/I_i)_{j-d_i}=\sum_i\shift^{d_i}(\vec s_i)(j) ,
\]
and as \((I_1,\dots,I_r)\) ranges over all tuples of monomial ideals --- unconstrained by one
another: \((\vec s_1,\dots,\vec s_r)\) ranges over all of \(\mathcal F_n^{\,r}\).

(b) Restricted to \textbf{lex} ideals, \(I\mapsto\HF(E/I)\) is a bijection onto \(\mathcal F_n\). It is well
defined (a lex ideal is monomial); injective, since a lex ideal is degreewise the lex-initial segment
of \(\mathrm{Mon}(E_j)\) of size \(\binom nj-s(j)\) and is therefore determined by \(\vec s\); and
surjective by Kruskal--Katona in its exterior form \textcite[Theorem 4.1]{AramovaHerzogHibi1997}, every \(f\)-vector being attained by a compressed
family, whose complementary ideal is lex.

(c) By Amata and Crupi's Definition 2.6 \(\Leftrightarrow\) 2.7, a lex submodule of \(F\) is exactly
an \(L=\bigoplus_iL_ig_i\) with every \(L_i\) a lex ideal and
\(\mathfrak m^{\rho_i+d_i-d_{i-1}}\subseteq L_{i-1}\) for \(i=2,\dots,r\), where
\(\rho_i=\indeg(L_i)\). The exponent is \(\ge0\) because \(\rho_i\ge0\) and \(d_i\ge d_{i-1}\), so
Lemma \ref{lem-dictionary} applies and turns the side condition into
\(\tp(\vec s_{i-1})<\indeg(\vec s_i)+\delta_i\). With (b) componentwise, \(L\mapsto(\vec s_1,\dots,\vec s_r)\) is a
bijection from lex submodules onto the chained tuples.

(d) \(L\mapsto\HF(F/L)\) is a bijection from lex submodules onto \(\mathrm{FF}(\vec d)\). Injective: by
Definition 2.6, \(\mathrm{Mon}_j(L)\) is the initial segment of \(\mathrm{Mon}_j(F)\) of size
\(\dim F_j-\HF(F/L)(j)\) in the module order, and an initial segment is determined by its
cardinality, so \(L\) is determined degreewise. Surjective: monomial submodules are graded, so
\(\mathrm{FF}(\vec d)\subseteq\lbrace\HF(F/M):M\ \text{graded}\rbrace=\lbrace\HF(F/L):L\ \text{lex}\rbrace\)
by Amata and Crupi's Theorem 4.2, implication (a) \(\Leftrightarrow\) (c); and conversely a lex
submodule is monomial, so the reverse inclusion holds by (a).

Composing (c) and (d) gives the claim.
\end{proof}

\begin{remark}
(\emph{An explicit inverse, and what it still needs.}) The proof above produces the chained representative
through Amata--Crupi's lex submodule, hence algebraically. A \textbf{closed-form} inverse is now known:
writing \(Q_i(j)=\max\bigl(0,\;x(j)-\sum_{k>i}\binom n{j-d_k}\bigr)\), the chained preimage of
\(\vec x\in\mathrm{FF}(\vec d)\) is the degreewise clamp
\begin{equation}\label{eq:water-filling}
  \hat s_i(m)=Q_i(m+d_i)-Q_{i-1}(m+d_i),
\end{equation}
a water-filling from the last component. That \eqref{eq:water-filling} is the \textbf{unique} chained tuple
summing to \(\vec x\) is proved using downward closure alone --- no Kruskal--Katona and no algebra ---
and the construction commutes with the involution of Remark \ref{rem-gap-reversal} on the nose, so the
symmetry that any bijective proof must respect is satisfied here by construction.

Two features of that half are worth recording, both of them visible only once it was formalised. It
needs no sortedness, provided the chaining condition is read in the absolute form
\eqref{eq:chaining-absolute} of Remark \ref{rem-sorted} below. And the \textbf{structure} it rests on --- that a chained tuple has pairwise
disjoint, increasing free windows \([\indeg(\vec s_i)+d_i,\;\tp(\vec s_i)+d_i]\), so that in any one
degree at most one component is free while the earlier ones have vanished and the later ones are still
full binomials --- uses only that condition together with
\(\indeg\le\tp+1\). In particular it holds for arbitrary non-negative integer vectors satisfying those
inequalities, whether or not they are \(f\)-vectors: being an \(f\)-vector is needed just for the
single step identifying a below-window entry as the full binomial \(\binom n{\cdot}\). The clamp is
thus elementary in a stronger sense than "no Kruskal--Katona": most of it is not about \(f\)-vectors
at all.

What is \textbf{not} yet combinatorial is the other half: that the clamps \eqref{eq:water-filling} are
\(f\)-vectors at all. That is true --- it is Theorem \ref{thm-chaining} again, through the lex
representative --- but proving it directly is open. For \(r=2\) it reduces to two properties of the
Kruskal--Katona shadow function, one of which we can prove and the other of which is verified far
beyond any range the argument needs but remains open on a sporadic set; for \(r\ge3\) a further
inequality is needed. So the bijection is explicit and one-directionally elementary, and open
problem 4 of Section \ref{sec-verification} stands.
\label{rem-explicit-inverse}
\end{remark}

\begin{remark}
Three hypotheses deserve to be stated rather than left implicit.

\textbf{The order convention.} Amata and Crupi's module order makes \(g_1\)'s monomials the largest, so a
lex submodule fills \(g_1\) first and the chaining condition constrains the \textbf{earlier} component
\(I_{i-1}\) using the \textbf{later} component's initial degree. The inequality in Theorem \ref{thm-chaining}
is oriented accordingly.

\textbf{Sortedness.} \(d_1\le\cdots\le d_r\) is required as stated, because it is what makes every exponent
\(\rho_i+\delta_i\) non-negative, which is the range in which Lemma \ref{lem-dictionary}(1) holds. This
costs nothing: \(\mathrm{FF}(\vec d)\) is symmetric in its summands, so an unsorted shift vector should be
sorted first. The chained count as written is simply not the right-hand side for unsorted \(\vec d\).

It is worth being precise about \textbf{where} sortedness is spent, because it is spent in one place only.
Write the chaining condition in the equivalent \textbf{absolute} form
\begin{equation}\label{eq:chaining-absolute}
  \tp(\vec s_{i-1})+d_{i-1}\;<\;\indeg(\vec s_i)+d_i \qquad(i=2,\dots,r),
\end{equation}
which is the condition of Theorem \ref{thm-chaining} rearranged. For sorted \(\vec d\) the two forms are literally
the same inequality. What sortedness buys is the passage to the \textbf{printed} form with the gap
\(\delta_i=d_i-d_{i-1}\) as a non-negative integer, i.e. Lemma \ref{lem-dictionary}(1); and it buys only
one of the two implications, since \eqref{eq:chaining-absolute} gives the printed inequality for every
\(\vec d\), sorted or not. Everything downstream that uses the chaining condition only as a condition
(the disjoint-window normal form and the uniqueness half of Remark
\ref{rem-explicit-inverse}) needs no sortedness at all when \eqref{eq:chaining-absolute} is taken as
the definition. This was established while formalising the argument (Section
\ref{sec-verification}), where the hypothesis turned out to be usable nowhere else.

\textbf{Which half of Theorem 4.2.} Only implication (a) \(\Leftrightarrow\) (c) is used --- existence and
uniqueness of a lex representative. That half is sound. The clause (b) of the same theorem is
\textbf{false as printed} for tied initial generator degrees, a defect treated in the third paper of this
series; nothing here depends on it.
\label{rem-sorted}
\end{remark}
\section{The counting theorem}
\label{sec-main}
\begin{definition}
(\emph{The initial vector and the transfer matrices.}) With \(A\) as in Definition \ref{def-A}, define a
\textbf{row vector} \(\vec u\) and, for each integer \(\delta\ge0\), a \textbf{square matrix} \(T(\delta)\) by
\begin{equation}\label{eq:def-uT}
  u[\tau]=\sum_{\rho}A[\tau][\rho],
  \qquad
  T(\delta)[\tau][\tau']=\sum_{\rho>\tau-\delta}A[\tau'][\rho].
\end{equation}
\textbf{Format, stated explicitly because both indices run over the same set and it is easy to mistake one
for the other.} Both are indexed by values of \(\tp\),
\begin{equation}\label{eq:def-uT-range}
  \tau,\tau'\in\lbrace-1,0,1,\dots,n\rbrace ,
\end{equation}
so \(\vec u\) has \(n+2\) entries and \(T(\delta)\) is \((n+2)\times(n+2)\). In \(T(\delta)\) \textbf{the row
index \(\tau\) is the \(\tp\) of the previous summand and the column index \(\tau'\) is the \(\tp\) of
the current one}; the summation variable \(\rho\) is the \(\indeg\) of the current summand and is
summed out, so it is not an index of \(T\). Thus \(u[\tau]\) counts the \(\vec s\in\mathcal F_n\)
with \(\tp(\vec s)=\tau\), and \(T(\delta)[\tau][\tau']\) counts those with \(\tp(\vec s)=\tau'\)
whose \(\indeg\) is large enough to be chained after a predecessor of \(\tp\) equal to \(\tau\).
Finally \(\vec 1\) is the column vector of \(n+2\) ones. Note that \(\vec u\) is the row sum of
\(A\), and that \(T(\delta)\) is \textbf{increasing entrywise in \(\delta\)}, since enlarging \(\delta\)
only enlarges the range of \(\rho\).
\label{def-uT}
\end{definition}

\begin{theorem}
Fix \(n\ge1\), \(r\ge1\) and integers \(d_1\le\cdots\le d_r\), and let \(\delta_i=d_i-d_{i-1}\) for
\(2\le i\le r\) be the gaps. With \(\vec u\), \(T(\delta)\) and \(\vec 1\) as in
Definition \ref{def-uT}, built from the matrix \(A\) of Definition \ref{def-A} (equivalently, by
Lemma \ref{lem-joint}, assembled from one run of Linusson's recursion at \(p=n\)), we have
\begin{equation}\label{eq:auto-15}
  \lvert\mathrm{FF}(\vec d)\rvert\;=\;\vec u\cdot T(\delta_2)\cdot T(\delta_3)\cdots T(\delta_r)\cdot\vec 1 .
\end{equation}
\label{thm-main}
\end{theorem}

\begin{proof}
By Theorem \ref{thm-chaining}, \(|\mathrm{FF}(\vec d)|\) counts tuples with
\(\tp(\vec s_{i-1})<\indeg(\vec s_i)+\delta_i\). Constraint \(i\) involves \(s_{i-1}\) only through \(\tp\) and
\(\vec s_i\) only through \(\tp\) and \(\indeg\). Summing over tuples in the order \(\vec s_1,\vec s_2,\dots,\vec s_r\)
and grouping each \(\vec s_i\) by its \((\tp,\indeg)\) class therefore gives the stated product: \(u[\tau]\)
counts the admissible \(\vec s_1\) with \(\tp=\tau\), and \(T(\delta_i)[\tau][\tau']\) counts the \(\vec s_i\)
with \(\tp=\tau'\) whose \(\indeg=\rho\) satisfies \(\tau<\rho+\delta_i\).
\end{proof}

\begin{mexample}
Continuing Example \ref{ex-A} at \(n=3\). The row sums of the matrix displayed there give the vector
\(\vec u=(1,1,3,4,1)\), and \eqref{eq:def-uT} then produces
\[
  T(0)=\begin{pmatrix}
    1&1&3&4&1\\ 0&1&3&4&1\\ 0&0&1&3&1\\ 0&0&0&1&1\\ 0&0&0&0&1
  \end{pmatrix},
  \quad
  T(1)=\begin{pmatrix}
    1&1&3&4&1\\ 1&1&3&4&1\\ 0&1&3&4&1\\ 0&0&1&3&1\\ 0&0&0&1&1
  \end{pmatrix},
  \quad
  T(4)=\begin{pmatrix}
    1&1&3&4&1\\ 1&1&3&4&1\\ 1&1&3&4&1\\ 1&1&3&4&1\\ 1&1&3&4&1
  \end{pmatrix}.
\]
Three things are visible. \(T(0)\) is upper triangular with unit diagonal, which is
Proposition \ref{prop-delta0}. Increasing \(\delta\) by one shifts the whole pattern down one row, since
the constraint \(\rho>\tau-\delta\) weakens by exactly one step --- this is the entrywise monotonicity
noted in Definition \ref{def-uT}. And at \(\delta=4=n+1\), the threshold of
Proposition \ref{prop-threshold}, every row has become \(\vec u\) itself: the matrix is
\(\vec 1\,\vec u\), of rank one, and its only non-zero eigenvalue is
\(\vec u\cdot\vec 1=10=\lvert\mathcal F_3\rvert\), which is Proposition \ref{prop-perron-threshold} in
this instance. Reading off \(r=2\), \(\vec u\,T(0)\,\vec 1=43\) and \(\vec u\,T(4)\,\vec 1=100=10^2\).
\label{ex-uT}
\end{mexample}

The rank enters only as the length of the product. One computation of \(A\) at fixed \(n\) answers
every \(r\) and every shift vector at once, which is what makes values far outside brute-force range
available: \(|\mathrm{FF}(8;0,0)|=27655902\) is immediate, while direct enumeration of the sumset
becomes infeasible around \(n=6\), \(r=2\).

\begin{corollary}
For \(r=1\), \(|\mathrm{FF}|=\sum_\tau u[\tau]=F^n(n)=|\mathcal F_n|\), which is Linusson's own count.
\label{cor-r1}
\end{corollary}

\begin{proposition}
\(|\mathrm{FF}(\vec d)|=|\mathcal F_n|^{\,r}\) if and only if every gap satisfies \(\delta_i\ge n+1\).
At \(\delta_i=n\) the constraint fails for exactly one pair, namely
\((\vec s_{i-1},\vec s_i)=(\text{full simplex},\text{void family})\).
\label{prop-threshold}
\end{proposition}

\begin{proof}
The constraint \(\tau<\rho+\delta\) is vacuous exactly when \(\delta>\max\tau-\min\rho\). Over
\(\mathcal F_n\), \(\max\tp=n\) (full simplex) and \(\min\indeg=0\) (void family), so the threshold is
\(\delta\ge n+1\); at \(\delta=n\) the single violating pair is the one realising both extremes.
\end{proof}

For \(r=2\) this gives \(|\mathcal F_n|^2-1\) at \(\delta=n\), i.e. \(99\), \(675\), \(9215\) for
\(n=3,4,5\), which matches the computed table.

\begin{proposition}
If all gaps vanish then \(T(0)\) is unipotent, meaning that \(T(0)-I\) is nilpotent, and
\(|\mathrm{FF}(n;0,\dots,0)|\) is a polynomial in
\(r\) of degree at most \(n+1\).
\label{prop-delta0}
\end{proposition}

\begin{proof}
\(T(0)[\tau][\tau']\ne0\) requires some \(\rho>\tau\) with \(A[\tau'][\rho]>0\); since
\(\rho\le\tau'+1\) always, this forces \(\tau'\ge\tau\), so \(T(0)\) is triangular in the
\(\tp\)-ordering. Its diagonal entries are \(T(0)[\tau][\tau]=A[\tau][\tau+1]=1\), the unique
\(s\) with \(\tp=\tau\) and \(\indeg=\tau+1\) being the full truncation at \(\tau\). So
\(N=T(0)-I\) is strictly triangular, hence nilpotent with \(N^{n+2}=0\); expanding
\((I+N)^{r-1}\) by the binomial theorem then terminates after \(n+2\) terms and is polynomial in
\(r\) of degree at most \(n+1\), the matrix size minus one. That expansion is
\eqref{eq:delta0-expansion}, where it is used again to pin the degree exactly.
\end{proof}

The bound of Proposition \ref{prop-delta0} is attained, and the leading coefficient is explicit. Both
follow from two structural entries of \(A\).

\begin{lemma}
For \(0\le j\le n\) one has \(A[j][j+1]=1\), the unique witness being the full truncation at \(j\);
and for \(1\le j\le n\) one has \(A[j][j]=\binom nj-1\). Moreover \(A[0][0]=0\), and row \(\tau=-1\)
of \(A\) is the single entry \(A[-1][0]=1\).
\label{lem-A-structure}
\end{lemma}

\begin{proof}
If \(\tp(\vec s)=j\) and \(\indeg(\vec s)=j+1\) then \(s(i)=\binom ni\) for \(i\le j\) and \(s(i)=0\)
for \(i>j\), so \(\vec s\) is the \(f\)-vector of the \(j\)-skeleton of the full simplex, and
conversely that vector has \(\tp=j\) and \(\indeg=j+1\); one witness exactly. If instead
\(\tp(\vec s)=\indeg(\vec s)=j\) then \(s(i)=\binom ni\) for \(i<j\), then
\(0<s(j)<\binom nj\) (positive because \(\tp=j\), deficient because \(\indeg=j\)), and
\(s(i)=0\) for \(i>j\). So \(\vec s\) is determined by \(m=s(j)\in\lbrace1,\dots,\binom nj-1\rbrace\),
and every such \(m\) occurs: adjoin any \(m\) of the \(j\)-subsets to the full \((j-1)\)-skeleton,
which stays an order ideal because every \((j-1)\)-subset is already present. The remaining two
statements are the void family and the two degenerate rows.
\end{proof}

\begin{theorem}
(\emph{The exact degrees at \(\delta=0\).}) Let \(n\ge1\). Then \(|\mathrm{FF}(n;0,\dots,0)|\) is a
polynomial in \(r\) of degree \textbf{exactly} \(n+1\), and \(|\mathrm{FF}_{\mathrm{pr}}(n;0,\dots,0)|\) one
of degree \textbf{exactly} \(n-1\), with leading coefficients
\begin{equation}\label{eq:delta0-leading}
  \frac{1}{(n+1)!}\prod_{j=0}^{n}\binom nj
  \qquad\text{and}\qquad
  \frac{1}{(n-1)!}\prod_{j=2}^{n}\binom nj
\end{equation}
respectively.
\label{thm-delta0-degrees}
\end{theorem}

\begin{proof}
Write \(N=T(0)-I\), strictly upper triangular in the \(\tp\)-ordering by Proposition
\ref{prop-delta0}. Expanding,
\begin{equation}\label{eq:delta0-expansion}
  \vec u\,T(0)^{r-1}\vec 1=\vec u\,(I+N)^{r-1}\vec 1
  =\sum_{k\ge0}\binom{r-1}{k}\,c_k,\qquad c_k=\vec u\,N^k\vec 1 ,
\end{equation}
and \(\binom{r-1}{k}\) is a polynomial in \(r\) of degree exactly \(k\). As the \(\binom{r-1}{k}\)
are linearly independent, \textbf{the degree is the largest \(k\) with \(c_k\ne0\)}, and the leading
coefficient is that \(c_k\) divided by \(k!\).

Since \(A\) is supported on \(\rho\le\tau'+1\), the superdiagonal of \(N\) is, by
Lemma \ref{lem-A-structure},
\begin{equation}\label{eq:delta0-superdiag}
  N[\tau][\tau+1]=A[\tau+1][\tau+1]+A[\tau+1][\tau+2]
  =\Bigl(\tbinom n{\tau+1}-1\Bigr)+1=\tbinom n{\tau+1}
  \qquad(-1\le\tau\le n-1),
\end{equation}
Pascal's row \(n\) exactly. For a strictly upper triangular matrix of size \(m\), the power
\(N^{m-1}\) has a single possibly non-zero entry, in the corner, equal to the product of the
superdiagonal --- only the full staircase path survives. Here \(m=n+2\), so
\(c_{n+1}=u[-1]\cdot\prod_{j=0}^{n}\binom nj\), and \(u[-1]=1\) because the void family is the unique
\(\vec s\) with \(\tp=-1\). This is non-zero, which gives the improper claim.

For the proper variant one adjustment is needed: the proper matrix is \textbf{not} unipotent on the full
index set, its columns \(\tau'=-1\) and \(\tau'=0\) vanishing identically. But \(\vec u\) vanishes on
those two coordinates as well, so every product \(\vec u\,T'(0)^k\) agrees with the corresponding
product over the principal block on \(\tau\in\lbrace1,\dots,n\rbrace\), and \textbf{that} block is
unipotent. On it the superdiagonal is again \eqref{eq:delta0-superdiag}, with column indices
\(\ge2\) so that the two zeroed columns are invisible, and \(u'[1]=A[1][2]=1\), the unique proper
vector with \(\tp=1\) being \((1,n,0,\dots,0)\). The block has size \(n\), so
\(c'_{n-1}=\prod_{j=2}^{n}\binom nj\ne0\).
\end{proof}

\begin{mexample}
The polynomials themselves, for \(n\le5\), obtained by interpolating
Theorem \ref{thm-main} at enough values of \(r\):
\begingroup\small
\begin{center}
\begin{tabular}{@{}cl@{}}
  \toprule
  \(n\) & \(\lvert\mathrm{FF}(n;0^r)\rvert\) \\
  \midrule
  1 & \(\tfrac12r^2+\tfrac32r+1\)\\
  2 & \(\tfrac13r^3+\tfrac32r^2+\tfrac{13}6r+1\)\\
  3 & \(\tfrac38r^4+\tfrac{23}{12}r^3+\tfrac{29}8r^2+\tfrac{37}{12}r+1\)\\
  4 & \(\tfrac45r^5+4r^4+\tfrac{47}6r^3+8r^2+\tfrac{131}{30}r+1\)\\
  5 & \(\tfrac{125}{36}r^6+\tfrac{185}{12}r^5+\tfrac{1915}{72}r^4+\tfrac{79}3r^3+\tfrac{1219}{72}r^2+\tfrac{25}4r+1\)\\
  \addlinespace
  \(n\) & \(\lvert\mathrm{FF}_{\mathrm{pr}}(n;0^r)\rvert\) \\
  \midrule
  1 & \(1\)\\
  2 & \(r+1\)\\
  3 & \(\tfrac32r^2+\tfrac52r+1\)\\
  4 & \(4r^3+7r^2+4r+1\)\\
  5 & \(\tfrac{125}6r^4+\tfrac{85}3r^3+\tfrac{79}6r^2+\tfrac{20}3r+1\)\\
  \bottomrule
\end{tabular}
\end{center}
\endgroup
Every one has constant term \(1\), as it must: at \(r=0\) the empty sum is the single point \(\vec0\).
The degrees and leading coefficients were checked against
\eqref{eq:delta0-leading} for \(n=1,\dots,7\) and agree in all fourteen cases; the two polynomials
at \(n=6,7\) are omitted here only for width, their leading coefficients being \(225/7\) and
\(84035/128\) improper, \(225\) and \(84035/16\) proper. One feature is worth flagging because it
is easy to assume otherwise: \textbf{the coefficients are not all positive}. The first negative one appears
at \(n=6\) in the proper polynomial and at \(n=7\) in the improper, so no positivity argument can be
read into Theorem \ref{thm-delta0-degrees}.
\label{tab-delta0-polys}
\end{mexample}

\begin{remark}
(\emph{Both degrees are longest paths.}) There is a reading of Theorem \ref{thm-delta0-degrees} under which
its two exponents are not computed at all but read off a picture, and it explains why the one step of
the proof that needed care is the one that did.

Let \(G\) be the weighted digraph on the vertex set \(\lbrace-1,0,\dots,n\rbrace\) of \(\tp\)-values
whose adjacency matrix is \(N=T(0)-I\). Unipotence of \(T(0)\) says exactly that \(G\) is a \textbf{directed
acyclic graph}: a non-negative matrix is unipotent precisely when every vertex carries a loop and
there are no other directed cycles, since all \(n+2\) eigenvalues being \(1\) forces
\(\operatorname{tr}(T(0)^k)=n+2\) for every \(k\), and the loops alone already supply that many closed
walks of each length. Removing the loops leaves \(N\), and a non-negative matrix is nilpotent exactly
when its digraph is acyclic. Here the arcs run \(\tau\to\tau'\) only for \(\tau'>\tau\), which is the
triangularity, and the arc \(\tau\to\tau+1\) carries weight \(\binom n{\tau+1}\) by
\eqref{eq:delta0-superdiag}.

Now \((N^k)[\tau][\tau']\) is the total weight of the \textbf{paths} of length exactly \(k\) from \(\tau\) to
\(\tau'\), paths and not walks because \(G\) is acyclic. So the expansion
\eqref{eq:delta0-expansion} reads
\begin{equation}\label{eq:delta0-paths}
  a_r=\sum_{k\ge0}\binom{r-1}k\cdot\bigl(\text{total weight of paths of length }k\text{ in }G\bigr),
\end{equation}
the paths being weighted by \(\vec u\) at their start and summed over their end. Since
\(\binom{r-1}k\) has degree exactly \(k\), this gives at once:

\begin{itemize}
\item the \textbf{\textbf{degree}} is the length of the longest path in \(G\) that \(\vec u\) reaches;
\item the \textbf{\textbf{leading coefficient}} is the total weight of those longest paths, divided by that length's
factorial.
\end{itemize}

Both halves of the theorem are now visible. In the improper case a path of length \(n+1\) must pass
through all \(n+2\) vertices, so it is the staircase \(-1\to0\to\cdots\to n\), unique, of weight
\(\prod_{j=0}^n\binom nj\): degree \(n+1\), leading coefficient
\(\prod_{j=0}^n\binom nj\big/(n+1)!\). In the proper case the vertices \(\tau=-1\) and \(\tau=0\) drop
out (\(\vec u'\) vanishes on them, and columns \(-1\) and \(0\) of \(T'(0)\) vanish too, so nothing
reaches them either), leaving a DAG on \(\lbrace1,\dots,n\rbrace\) whose longest path
\(1\to2\to\cdots\to n\) has length \(n-1\) and weight \(\prod_{j=2}^n\binom nj\).

This also says which step of the proof carries the content. "Only the full staircase survives", the
corner evaluation, is the assertion that \(G\) has exactly one longest path; everything else is the
binomial expansion. And it explains the shape of the answer: a product of binomial coefficients
appears because a longest path is forced to take every arc of the staircase, and a factorial divides
it because \(\binom{r-1}{k}\) does.
\label{rem-delta0-digraph}
\end{remark}

\begin{remark}
The improper leading coefficient is \(\prod_{j=0}^{n}\binom nj\big/(n+1)!\), the numerator being the
product of a whole row of Pascal's triangle (OEIS A001142). At \(n=2\) the two coefficients are
\(2/3!=1/3\) and \(1/1!=1\), matching the improper values \(14,30,55,91\) for \(r=2,\dots,5\), a
cubic with third difference \(2\), and the proper values
\(|\mathrm{FF}_{\mathrm{pr}}(2;0^r)|=r+1\). Theorem \ref{thm-delta0-degrees} was verified against
independent computation of both counts for \(n=2,\dots,8\) and \(r\) up to \(n+4\); see
Section \ref{sec-verification}.
\label{rem-delta0-degrees}
\end{remark}
\section{The proper variant}
\label{sec-proper}
The companion paper's polytope is built from \(S_n\), not \(\mathcal F_n\): its Hilbert functions come
from \textbf{proper} monomial submodules. The count restricts accordingly, but the restriction is not free:
it costs work in the proof below, and it costs a symmetry outright (Remark
\ref{rem-proper-not-reversible}).
The argument of Theorem \ref{thm-chaining} goes through the lex representative, so what must be checked
is that the lex representative of a properly achievable Hilbert function is itself proper.

\begin{lemma}
Let \(\vec x=\sum_i\shift^{d_i}(\vec s_i)\) with every \(\vec s_i\in S_n\). Then the unique lex submodule
\(L\) with \(\HF(F/L)=\vec x\) is proper.
\label{lem-lexproper}
\end{lemma}

\begin{proof}
By Definition 2.6, \(\mathrm{Mon}_j(L)\) is the initial segment of \(\mathrm{Mon}_j(F)\) of size
\(D_j=\dim F_j-x(j)\), and the segment fills the slots \(g_1,g_2,\dots\) in order. So \(L\) is proper
exactly when, for every \(i\), the segment does not reach \(g_i\)'s slot in the two module degrees
\(j=d_i\) and \(j=d_i+1\), that is, when \(D_j\le\sum_{k<i}\binom n{j-d_k}\) there. Writing
\(\binom nm=0\) and \(s_k(m)=0\) outside \(0\le m\le n\),
\[
  D_j=\sum_k\Bigl(\tbinom n{j-d_k}-s_k(j-d_k)\Bigr).
\]
For \(k\ge i\) and \(j\le d_i+1\) we have \(j-d_k\le d_i+1-d_k\le1\), and a proper \(\vec s_k\) has
\(s_k(0)=1=\binom n0\) and \(s_k(1)=n=\binom n1\), while both vanish in negative degrees; so every
such summand is zero. Hence \(D_j=\sum_{k<i}\bigl(\binom n{j-d_k}-s_k(j-d_k)\bigr)\le
\sum_{k<i}\binom n{j-d_k}\), which is exactly the capacity of the slots preceding \(g_i\).
\end{proof}

\begin{theorem}
Fix \(n\ge1\), \(r\ge1\) and a sorted \(\vec d\). Let \(A'\) be \(A\) with the columns \(\rho=0\) and
\(\rho=1\) set to zero, \(\vec u'\) its row sums, and \(T'(\delta)\) built from \(A'\) as in
Theorem \ref{thm-main}. Then
\begin{equation}\label{eq:auto-16}
  \lvert\mathrm{FF}_{\mathrm{pr}}(\vec d)\rvert
  =\#\bigl\lbrace(\vec s_1,\dots,\vec s_r)\in S_n^{\,r}:\tp(\vec s_{i-1})<\indeg(\vec s_i)+\delta_i\bigr\rbrace
  = \vec u'\cdot T'(\delta_2)\cdots T'(\delta_r)\cdot\vec 1 .
\end{equation}
\label{thm-proper}
\end{theorem}

\begin{proof}
\(S_n=\lbrace \vec s\in\mathcal F_n:\indeg(\vec s)\ge2\rbrace\), which is what deleting the two columns
implements. As in Theorem \ref{thm-chaining}(d), \(L\mapsto\HF(F/L)\) is injective on lex submodules; it
carries proper lex submodules into \(\mathrm{FF}_{\mathrm{pr}}(\vec d)\), since a proper lex submodule is a
proper monomial submodule, and onto it: given \(\vec x\in\mathrm{FF}_{\mathrm{pr}}(\vec d)\), Theorem 4.2
(a) \(\Leftrightarrow\) (c) supplies the lex representative and Lemma \ref{lem-lexproper} shows it is
proper. Definition 2.6 \(\Leftrightarrow\) 2.7 with Lemma \ref{lem-dictionary}, restricted componentwise
to lex ideals of initial degree \(\ge2\), identifies proper lex submodules with \(S_n\)-chained
tuples. The factorisation is then the rearrangement of Theorem \ref{thm-main} verbatim.
\end{proof}

\begin{proposition}
For \(n\ge1\), \(|\mathrm{FF}_{\mathrm{pr}}(\vec d)|=|S_n|^{\,r}\) if and only if every gap satisfies
\(\delta_i\ge n-1\).
\label{prop-proper-threshold}
\end{proposition}

\begin{proof}
As in Proposition \ref{prop-threshold}, with \(\max\tp=n\) and \(\min\indeg=2\) over \(S_n\).
\end{proof}

\begin{remark}
Two sequences recorded in earlier work on this project are reproduced exactly by
Theorem \ref{thm-proper} and extended. For \(r=2\), \(n=2,\dots,9\):
\begin{equation}\label{eq:auto-17}
  \delta=0:\quad 3,\ 12,\ 69,\ 627,\ 9904,\ 296289,\ 18009389,\ 2346218172,
\end{equation}
\begin{equation}\label{eq:auto-18}
  \delta=1:\quad 4,\ 22,\ 162,\ 1858,\ 38232,\ 1571774,\ 138881518,\ 27635025278 .
\end{equation}
The prefixes through \(n=6\) are the independently enumerated ones. Neither sequence, nor \(|S_n|\)
itself, is currently in the OEIS.
\label{rem-proper-sequences}
\end{remark}
\section{Growth}
\label{sec-growth}
\begin{proposition}
For a constant gap \(\delta\), \(a_r:=|\mathrm{FF}(n;0,\delta,2\delta,\dots)|=\vec u\cdot
T(\delta)^{r-1}\cdot\vec 1\) satisfies the linear recurrence whose characteristic polynomial is
that of \(T(\delta)\), and its exponential growth rate is the Perron root \(\lambda\) of
\(T(\delta)\), in the sense that \(a_r^{1/r}\to\lambda\). The same holds for the proper count with
\(T'(\delta)\).
\label{prop-recurrence}
\end{proposition}

\begin{proof}
Cayley--Hamilton applied to \(T(\delta)\), the vectors \(\vec u\) and \(\vec 1\) being fixed; and for
the growth rate, the spectral radius of a matrix is the exponential growth rate of the norms of its
powers, with \(\lambda\) that radius by Perron--Frobenius, \(T(\delta)\) being non-negative.

It is worth saying what this does \textbf{not} claim, since the natural stronger reading is false. There is
in general no constant \(C\) with \(a_r\sim C\lambda^r\): at \(\delta=0\) the matrix is unipotent, so
\(\lambda=1\), while \(a_r\) is a polynomial in \(r\) of degree exactly \(n+1\) by
Theorem \ref{thm-delta0-degrees} and \(a_r/\lambda^r\) is therefore unbounded. Such a \(C\) exists when
the dominant eigenvalue is simple, which the table below is consistent with for \(\delta\ge1\), but we
do not prove it.
\end{proof}

\begin{mexample}
At \(n=2\), \(\delta=1\), the characteristic polynomial of \(T(1)\) is
\(\lambda(\lambda-1)(\lambda^2-4\lambda+1)\), so \(a_r=5a_{r-1}-5a_{r-2}+a_{r-3}\); the values are
\(5,20,76,285,1065\) for \(r=1,\dots,5\), and indeed \(5\cdot285-5\cdot76+20=1065\). The Perron root
is \(2+\sqrt3\approx3.732\). At \(n=3\), \(\delta=1\), it is
\(\lambda(\lambda^4-9\lambda^3+21\lambda^2-15\lambda+3)\), giving
\(a_r=9a_{r-1}-21a_{r-2}+15a_{r-3}-3a_{r-4}\), checked against \(10,69,426,2532,14847\).

In both cases a factor of \(\lambda\) can be dropped before reading off the recurrence, and it is
worth saying why rather than leaving it to be noticed. \(T(\delta)\) is singular for
\(\delta\ge1\): rows \(\tau=-1\) and \(\tau=0\) both equal \(\vec u\), since
\(\rho>-1-\delta\) and \(\rho>-\delta\) admit the same \(\rho\ge0\). A root at \(0\)
contributes a term \(c\cdot0^{\,r}\), which vanishes for \(r\ge1\), so the quotient annihilates
the sequence just as the characteristic polynomial does and gives a shorter recurrence. The proper
matrix shows the same phenomenon doubled: \(T'(0)\) has characteristic polynomial
\(\lambda^2(\lambda-1)^n\), the two factors of \(\lambda\) being the two identically zero
columns that Section \ref{sec-proper}'s invariant-block argument turns on.

Across \(n\) and \(\delta\) the Perron roots are as follows, with \(\lvert\mathcal F_n\rvert\) and
\(\lvert S_n\rvert\) in the last two columns for comparison; an entry in bold has reached its limit
exactly, which is Proposition \ref{prop-perron-threshold}.
\begingroup\small
\begin{center}
\begin{tabular}{@{}cc rr cc@{}}
  \toprule
  \(n\) & \(\delta\) & \(\lambda(\delta)\) & \(\lambda_{\mathrm{pr}}(\delta)\)
        & \(\lvert\mathcal F_n\rvert\) & \(\lvert S_n\rvert\)\\
  \midrule
  2 & 0 & 1 & 1 & 5 & 2\\
  2 & 1 & 3.7321 & \textbf{2} & &\\
  2 & 2 & 4.7913 & \textbf{2} & &\\
  2 & 3 & \textbf{5} & \textbf{2} & &\\
  \addlinespace
  3 & 1 & 5.8188 & 4.3028 & 10 & 5\\
  3 & 2 & 9.0138 & \textbf{5} & &\\
  3 & 3 & 9.8990 & \textbf{5} & &\\
  3 & 4 & \textbf{10} & \textbf{5} & &\\
  \addlinespace
  4 & 1 & 9.8557 & 8.7566 & 26 & 16\\
  4 & 2 & 20.5859 & 15.3485 & &\\
  4 & 3 & 25.2491 & \textbf{16} & &\\
  4 & 4 & 25.9615 & \textbf{16} & &\\
  4 & 5 & \textbf{26} & \textbf{16} & &\\
  \addlinespace
  5 & 1 & 17.7573 & 17.1251 & 96 & 70\\
  5 & 2 & 57.8411 & 52.1804 & &\\
  5 & 3 & 87.4360 & 69.3802 & &\\
  5 & 4 & 95.4659 & \textbf{70} & &\\
  5 & 5 & 95.9896 & \textbf{70} & &\\
  5 & 6 & \textbf{96} & \textbf{70} & &\\
  \bottomrule
\end{tabular}
\end{center}
\endgroup
At \(\delta=0\) both roots are \(1\), the unipotent case of Proposition \ref{prop-delta0}, where the
growth is polynomial rather than exponential. The proper root reaches its limit two steps before the
improper one, and the last improper step is always the narrowest.
\label{mex-recurrences}
\end{mexample}

\begin{proposition}
(\emph{Normal form: all the gap-dependence sits in one shift.}) Let \(P\) be the
\((n+2)\times(n+2)\) saturating one-step shift on the \(\tp\)-index,
\(P[\tau][\tau']=1\) if \(\tau'=\max(\tau-1,-1)\) and \(0\) otherwise. Write \(B=T(0)\). Then for
every \(\delta\ge0\),
\begin{equation}\label{eq:normalform}
  T(\delta)=P^{\delta}B,
  \qquad
  \vec u=e_{-1}^{\!\top}B,
  \qquad\text{so}\qquad
  \lvert\mathrm{FF}(\vec d)\rvert
   =e_{-1}^{\!\top}B\,P^{\delta_2}B\,P^{\delta_3}\cdots P^{\delta_r}B\,\vec 1 .
\end{equation}
\label{prop-normalform}
\end{proposition}

\begin{proof}
By Definition \ref{def-uT}, \(T(\delta)[\tau][\tau']=\sum_{\rho>\tau-\delta}A[\tau'][\rho]\) depends on
\(\tau\) and \(\delta\) only through \(\tau-\delta\), and coincides with row \(\tau-\delta\) of
\(T(0)\) whenever \(\tau-\delta\ge-1\). For \(\tau-\delta<-1\) the condition \(\rho>\tau-\delta\) is
vacuous, as is the condition \(\rho>-1\) defining row \(-1\) of \(T(0)\), so the row saturates at row
\(-1\). Shifting each row down by one and saturating at \(-1\) is exactly left multiplication by
\(P\). The second identity is row \(-1\) of \(T(0)\): \(T(0)[-1][\tau']=\sum_{\rho\ge0}A[\tau'][\rho]
=u[\tau']\).
\end{proof}

\begin{remark}
Two consequences are worth naming. First, \textbf{the count depends on each gap only through
\(\min(\delta_i,n+1)\)}, because \(P\) saturates: \(P^{\delta}=P^{n+1}=\vec1\,e_{-1}^{\!\top}\) for
every \(\delta\ge n+1\). Second, that same identity gives a second and shorter proof of
Propositions \ref{prop-threshold} and \ref{prop-perron-threshold}: above the threshold
\(T(\delta)=\vec1\,e_{-1}^{\!\top}B=\vec1\,\vec u\), manifestly of rank one with non-zero eigenvalue
\(\vec u\cdot\vec1=\lvert\mathcal F_n\rvert\). The pair \((B,P)\), one unipotent
matrix carrying the Pascal superdiagonal of Lemma \ref{lem-A-structure} and one \(0/1\) shift, encodes
every gap vector at once.

It also explains why the matrices do not commute while the improper count is nevertheless
reversal-symmetric
(Example \ref{ex-factors} and Remark \ref{rem-gap-reversal}): the non-commutation is visible in
\eqref{eq:normalform}, since \(P\) and \(B\) do not commute, whereas the reversal symmetry is the
persymmetry of \(A\) applied factor by factor inside the bilinear form, a property the individual
\(T(\delta)\) cannot express once \(\rho\) has been summed out.
\label{rem-normalform}
\end{remark}

\begin{proposition}
(\emph{The Perron root above the threshold.}) Let \(\lambda(\delta)\) be the Perron root of \(T(\delta)\)
and \(\lambda_{\mathrm{pr}}(\delta)\) that of its proper analogue. Then
\begin{equation}\label{eq:perron-threshold}
  \lambda(\delta)=\lvert\mathcal F_n\rvert\ \ \text{for every }\delta\ge n+1,
  \qquad
  \lambda_{\mathrm{pr}}(\delta)=\lvert S_n\rvert\ \ \text{for every }\delta\ge n-1 .
\end{equation}
In both cases the value is attained, not merely approached, and the thresholds are exactly those of
Propositions \ref{prop-threshold} and \ref{prop-proper-threshold}.
\label{prop-perron-threshold}
\end{proposition}

\begin{proof}
By Theorem \ref{thm-main}, \(T(\delta)[\tau][\tau']=\sum_{\rho>\tau-\delta}A[\tau'][\rho]\). The
condition \(\rho>\tau-\delta\) is vacuous exactly when \(\delta>\max\tp-\min\indeg\), which over
\(\mathcal F_n\) is \(\delta\ge n+1\), the computation already made in
Proposition \ref{prop-threshold}. When it is vacuous the inner sum no longer depends on \(\tau\):
every row of \(T(\delta)\) is the same vector \(c\), where \(c[\tau']=\sum_\rho A[\tau'][\rho]\) is
the number of \(\vec s\in\mathcal F_n\) with \(\tp(\vec s)=\tau'\). So \(T(\delta)=\vec1\,c^{\!\top}\)
has rank one, its non-zero eigenvalue is \(c^{\!\top}\vec1=\sum_{\tau'}c[\tau']=|\mathcal F_n|\), and
every other eigenvalue is \(0\). The proper statement is the same computation with \(\min\indeg=2\)
over \(S_n\), giving the threshold \(\delta\ge n-1\) and the root \(|S_n|\).
\end{proof}

\begin{remark}
This is the spectral shadow of Propositions \ref{prop-threshold} and \ref{prop-proper-threshold}, and it
is worth seeing why the two statements are the same fact twice. Those propositions say that above
the threshold the chaining constraint is vacuous, so the sumset is a full Cartesian product and
\(|\mathrm{FF}|\) is \(|\mathcal F_n|^{\,r}\) exactly; Proposition \ref{prop-perron-threshold} says the
same thing one level down, in the matrix: the constraint is what couples consecutive summands, and
with it gone \(T(\delta)\) forgets its row index entirely. A rank-one matrix \(\vec1\,c^{\!\top}\)
has a single non-zero eigenvalue, and it is the total count; multiplying \(r-1\) copies of it
reproduces \(|\mathcal F_n|^{\,r}\) directly.

Figure \ref{fig:perron-vs-delta} shows both roots at \(n=5\). The shape is the content of
Proposition \ref{prop-perron-threshold}: each climbs and then stops dead, the proper one two steps
earlier than the improper, and after that they are constant rather than asymptotic.

\begin{figure}[htbp]
\centering
\includegraphics[width=0.86\linewidth]{../img/perron-vs-delta-n5.pdf}
\caption{The Perron roots \(\lambda(\delta)\) and \(\lambda_{\mathrm{pr}}(\delta)\) of
Definition~\ref{def-uT}'s transfer matrices at \(n=5\), against the gap \(\delta\). Dotted lines mark
\(\lvert\mathcal F_5\rvert=96\) and \(\lvert S_5\rvert=70\); the vertical rules mark the two
thresholds of Proposition~\ref{prop-perron-threshold}, \(\delta=n-1\) for the proper root and
\(\delta=n+1\) for the improper. Each root is *equal* to its limit from its own threshold onward, not
merely close to it.}
\label{fig:perron-vs-delta}
\end{figure}

Below the threshold the root is strictly smaller, and it climbs. For \(n=3\), where
\(|\mathcal F_3|=10\) and \(|S_3|=5\), the improper roots at \(\delta=1,2,3,4\) are
\begin{equation}\label{eq:auto-19}
  5.8188\ldots,\quad 9.0138\ldots,\quad 9.8990\ldots,\quad 10,
\end{equation}
reaching \(|\mathcal F_3|\) exactly at \(\delta=4=n+1\) and staying there; the proper roots at
\(\delta=1,2\) are \(4.3028\ldots\) and \(5\), reaching \(|S_3|\) exactly at \(\delta=2=n-1\). The
approach is from below in both cases, as it must be, since \(T(\delta)\) increases entrywise with
\(\delta\). Note how narrow the last step is: at \(\delta=n\) the improper root is already
\(9.8990\ldots\) against \(10\), and at \(n=6\) it is \(552.998\ldots\) against \(553\), which is
Proposition \ref{prop-threshold}'s "the constraint fails for exactly one pair" made quantitative.
\label{rem-perron}
\end{remark}

The remaining growth question is in the number of vertices. The count \(|\mathrm{FF}(\vec d)|\) itself
cannot be polynomial in \(n\), since it dominates \(|\mathcal F_n|=F^n(n)\). The right generalisation
is Linusson's own: keep the dimension bound and free the vertex bound.

\begin{definition}
For \(p\ge n\), let \(\mathcal F_n^{(p)}\) be the set of \(f\)-vectors with \(f_1\le p\) and
\(f_j=0\) for \(j>n\), and \(\mathrm{FF}^{(p)}(n;\vec d)\) the corresponding sumset.
\label{def-twoparam}
\end{definition}

\begin{theorem}
Fix \(n\ge1\) and \(r\ge1\); let \(d_1\le\cdots\le d_r\) be integers and put
\(\delta_i=d_i-d_{i-1}\) for \(2\le i\le r\), so that \(\vec d\) is sorted and the gaps are
non-negative. Let \(p\ge n\), let \(\mathcal F_n^{(p)}\) be as in Definition \ref{def-twoparam}, and set
\begin{equation}\label{eq:Ap}
  A^{(p)}[\tau][\rho]=E^p(\tau,\rho-1)-E^p(\tau-1,\rho-1),
  \qquad \tau=-1,\dots,n,\quad\rho=0,\dots,n+1,
\end{equation}
with \(\vec u^{(p)}\) its vector of row sums and \(T^{(p)}(\delta)\) built from \(A^{(p)}\) by
\eqref{eq:def-uT}, both indexed as in Definition \ref{def-uT}. Then
\begin{equation}\label{eq:auto-20}
  \lvert\mathrm{FF}^{(p)}(n;\vec d)\rvert\;=\;\vec u^{(p)}\cdot T^{(p)}(\delta_2)\cdots T^{(p)}(\delta_r)\cdot\vec 1 ,
\end{equation}
and for fixed \(n\), \(r\) and \(\vec d\) this is a polynomial in \(p\).
\label{thm-poly-p}
\end{theorem}

\begin{proof}
Lemma \ref{lem-joint} and Theorems \ref{thm-chaining} and \ref{thm-main} go through verbatim over the exterior
algebra on \(p\) generators, with one extra ingredient: the lex representative must respect the
truncation, i.e. \(\mathfrak m^{n+1}F\subseteq M\) must force \(L\supseteq\mathfrak m^{n+1}F\). This
is the mirror image of the capacity argument of Lemma \ref{lem-lexproper}. In module degree
\(j=d_i+m\) with \(m\ge n+1\), every component \(k\le i\) has \(s_k(j-d_k)=0\) because
\(\tp(s_k)\le n\); hence \(D_j\ge\sum_{k\le i}\binom p{j-d_k}\) and the initial segment covers
\(g_i\)'s slot. Polynomiality in \(p\) is then immediate from Theorem \ref{thm-main}'s formula, each
entry \(E^p(\tau,\rho-1)\) being a polynomial in \(p\) for fixed \(\tau,\rho\) by Linusson's
Theorem 5.3.
\end{proof}

\textbf{Status.} The reduction above is proved; the polynomiality of the entries is Linusson's Theorem 5.3,
which we cite rather than reprove, so Theorem \ref{thm-poly-p} is proved modulo that citation. It was
verified against brute-force truncated sumsets for \(n\le3\), \(p\le5\), \(r\le3\), with no
mismatches, and the empirical degrees in \(p\) at \(r=2\) are \(3,5,6,9\) for
\((n,\delta)=(2,0),(2,1),(3,0),(3,1)\).

\begin{mexample}
The factors, written out. Take \(n=3\), \(r=3\) and \(\vec d=(0,1,3)\), so the gaps are
\((\delta_2,\delta_3)=(1,2)\) and the product of Theorem \ref{thm-main} is \(\vec u\,T(1)\,T(2)\,\vec1\)
with the matrices of Example \ref{ex-uT}. Multiplying left to right,
\[
  \vec u=(1,1,3,4,1)
  \ \xrightarrow{\ T(1)\ }\ (2,5,19,33,10)
  \ \xrightarrow{\ T(2)\ }\ (26,59,187,266,69),
\]
whose entries sum to \(\lvert\mathrm{FF}(3;0,1,3)\rvert=607\). Each application of a transfer matrix
advances one summand: the vector after \(k\) factors records, for each value of \(\tp\), how many
chained \((k+1)\)-tuples end in a summand with that \(\tp\).

Now reverse the gaps, \(\vec d=(0,2,3)\) and \((\delta_2,\delta_3)=(2,1)\). The intermediate vector is
different, \(\vec u\,T(2)=(5,9,28,39,10)\), but the total is again \(607\), which is the
gap-reversal symmetry of Remark \ref{rem-gap-reversal}. It is worth seeing that this is \textbf{not} because
the factors commute: at \(n=3\),
\[
  T(1)T(2)\ \ne\ T(2)T(1),
\]
already in the first row, \((5,9,28,39,10)\) against \((2,5,19,33,10)\). The symmetry is a property of
the bilinear form \(\vec u\,(\cdot)\,\vec1\) and of the duality behind it, not of the matrices.
\label{ex-factors}
\end{mexample}

Linusson's Theorem 5.3 and Proposition 5.4 give the degree and the leading coefficient of his own
count as a polynomial in \(p\). They transport through the product of Theorem \ref{thm-poly-p}, and the
result is clean. Write \(C_2(x)=\binom x2\) and, for each gap, \(\kappa_i=\max(0,n-\delta_i)\).

\begin{theorem}
(\emph{Degree and leading coefficient in \(p\).}) Fix \(n\ge1\), \(r\ge1\) and a sorted \(\vec d\) with
gaps \(\delta_i=d_i-d_{i-1}\). Then
\begin{equation}\label{eq:degree-in-p}
  \deg_p\lvert\mathrm{FF}^{(p)}(n;\vec d)\rvert
   \;=\; r\,C_2(n+1)\;-\;\sum_{i=2}^{r}C_2(\kappa_i+1),
\end{equation}
which is strictly increasing in each \(\delta_i\) up to \(\delta_i=n\) and constant thereafter. If
moreover every gap satisfies \(\delta_i\ge1\), then the summands decouple to leading order:
\begin{equation}\label{eq:lc-in-p}
  \lvert\mathrm{FF}^{(p)}(n;\vec d)\rvert\;\sim\;F^p(n)\prod_{i=2}^{r}E^p(n,\kappa_i)
  \qquad(p\to\infty).
\end{equation}
\label{thm-asymptotics-p}
\end{theorem}

\begin{proof}
By Lemma \ref{lem-joint} the entries of \(A^{(p)}\) are differences \(E^p(\tau,\rho-1)-E^p(\tau-1,\rho-1)\),
and Linusson's Theorem 5.3 gives \(\deg_pE^p(m,k)=C_2(m+1)-C_2(k+1)\), so the subtrahend has degree
strictly smaller (by exactly \(\tau\)) and cannot cancel the leading term. Hence
\(\deg_pA^{(p)}[\tau][\rho]=C_2(\tau+1)-C_2(\rho)\), with leading coefficient the positive quantity
of his Proposition 5.4.

Expanding the matrix product of Theorem \ref{thm-poly-p} as a sum over paths
\(\tau_1,\dots,\tau_r\), every term is a product of such entries, all with positive leading
coefficients, so again nothing cancels and the degree is the maximum over paths. The path degree
\textbf{separates position by position},
\begin{equation}\label{eq:path-degree}
  \sum_{j}h_{\delta_{j+1}}(\tau_j)+C_2(\tau_r+1),
  \qquad h_\delta(t)=C_2(t+1)-C_2\bigl(\max(1,t-\delta+1)\bigr),
\end{equation}
so maximising is not an optimisation over paths at all but a coordinatewise one. For \(\delta\ge1\)
the function \(h_\delta\) is strictly increasing on \(\lbrace0,\dots,n\rbrace\), so the maximiser is
unique (every \(\tau_j=n\)), which gives \eqref{eq:degree-in-p} and, since a unique maximising
path contributes a single product of leading coefficients, the product form \eqref{eq:lc-in-p}.
\end{proof}

\begin{remark}
At \(\delta_i=0\) the function \(h_0\) of \eqref{eq:path-degree} is identically zero, so the
maximiser is no longer unique and the leading coefficient becomes a \textbf{sum} over the degree-maximising
paths --- pinned to \(\tau=n\) wherever a positive gap follows, and free non-decreasing across runs
of zero gaps. The degree formula \eqref{eq:degree-in-p} is unaffected. The product form
\eqref{eq:lc-in-p} genuinely fails: at \(n=2\) with the single gap \(\delta_2=0\) the leading
coefficient is \(1\), against \(1/6\) for the product. This is the same phenomenon as
Proposition \ref{prop-delta0}, where vanishing gaps collapse the transfer matrix to a unipotent one,
seen now in the \(p\)-direction rather than the \(r\)-direction.

Two smaller observations. The degree and the leading coefficient are already those of the saturated
case at \(\delta_i=n\), one step before the threshold \(\delta_i\ge n+1\) of
Proposition \ref{prop-threshold}; the two regimes first differ in lower-order terms, the single
violating pair recorded there surviving at every \(p\). And the proper variant transports as
well --- Theorem \ref{thm-asymptotics-p-proper} below.
\label{rem-asymptotics-zero-gaps}
\end{remark}

The transport of the proper count is the same argument run on the principal block, and the outcome
is a uniform shift by one. Let \(S_n^{(p)}=\lbrace\vec s\in\mathcal F_n^{(p)}:\indeg(\vec s)\ge2\rbrace\),
of cardinality
\(\lvert S_n^{(p)}\rvert=\sum_{k=1}^{n}E^p(n,k)=F^p(n)-F^{p-1}(n)\), let
\(\mathrm{FF}^{(p)}_{\mathrm{pr}}(n;\vec d)\) be the corresponding sumset, and write
\(\kappa_i'=\max(1,n-\delta_i)\).

\begin{theorem}
(\emph{The proper analogue.}) Let \(A'^{(p)}\) be \(A^{(p)}\) with the columns \(\rho=0\) and
\(\rho=1\) set to zero, \(\vec u'^{(p)}\) its row sums, and \(T'^{(p)}(\delta)\) built from
\(A'^{(p)}\) by \eqref{eq:def-uT}. Then, for \(p\ge n\),
\begin{equation}\label{eq:proper-p-product}
  \lvert\mathrm{FF}^{(p)}_{\mathrm{pr}}(n;\vec d)\rvert
  \;=\;\vec u'^{(p)}\cdot T'^{(p)}(\delta_2)\cdots T'^{(p)}(\delta_r)\cdot\vec 1 ,
\end{equation}
a polynomial in \(p\) of degree
\begin{equation}\label{eq:proper-degree-in-p}
  \deg_p\lvert\mathrm{FF}^{(p)}_{\mathrm{pr}}(n;\vec d)\rvert
   \;=\; r\,C_2(n+1)\;-\;1\;-\;\sum_{i=2}^{r}C_2(\kappa_i'+1),
\end{equation}
strictly increasing in each \(\delta_i\) up to \(\delta_i=n-1\) and constant thereafter. If every
gap satisfies \(\delta_i\ge1\), then
\begin{equation}\label{eq:proper-lc-in-p}
  \lvert\mathrm{FF}^{(p)}_{\mathrm{pr}}(n;\vec d)\rvert\;\sim\;E^p(n,1)\prod_{i=2}^{r}E^p(n,\kappa_i')
  \qquad(p\to\infty),
\end{equation}
where \(E^p(n,1)\) may be replaced by \(\lvert S_n^{(p)}\rvert\), of which it is the dominant term.
\label{thm-asymptotics-p-proper}
\end{theorem}

\begin{proof}
The formula: the proofs of Theorems \ref{thm-proper} and \ref{thm-poly-p} compose. Over \(p\)
generators the capacity argument of Lemma \ref{lem-lexproper} reads \(s_k(0)=1=\binom p0\) and
\(s_k(1)=p=\binom p1\) and goes through verbatim, so the lex representative of a properly
achievable Hilbert function is again proper, and deleting the \(\indeg\le1\) columns implements
the restriction, as in Theorem \ref{thm-proper}.

The asymptotics: row \(\tau=-1\) of \(A^{(p)}\) is \(e_0\) and row \(\tau=0\) is \(e_1\), so both
die with the deleted columns. Hence \(\vec u'^{(p)}\) vanishes at \(\tau=-1,0\) and the columns
\(\tau'=-1,0\) of every \(T'^{(p)}(\delta)\) vanish, confining the whole product to the principal
block \(\tau\in\lbrace1,\dots,n\rbrace\), the restriction that already repaired the proper half of
Theorem \ref{thm-delta0-degrees}, now at every gap simultaneously. On the block every surviving
entry is an untouched entry of \(A^{(p)}\), with the degrees and positive leading coefficients of
Theorem \ref{thm-asymptotics-p}'s proof; the inner sums now start at \(\max(2,\tau-\delta+1)\)
rather than \(\max(1,\tau-\delta+1)\), so nothing cancels, and the path degree separates as in
\eqref{eq:path-degree} with \(h_\delta\) replaced by
\(h'_\delta(t)=C_2(t+1)-C_2(\max(2,t-\delta+1))\) and a path-independent constant \(-C_2(2)=-1\)
contributed by the first factor. On \(\lbrace1,\dots,n\rbrace\), \(h'_0\equiv0\) and \(h'_\delta\)
is strictly increasing for \(\delta\ge1\), so the unique maximiser is again the all-\(n\) path;
its factors have degrees \(C_2(n+1)-1\) and \(C_2(n+1)-C_2(\kappa_i'+1)\), with leading
coefficients those of \(E^p(n,1)\) and \(E^p(n,\kappa_i')\) by Linusson's Proposition 5.4, which
gives \eqref{eq:proper-degree-in-p} and \eqref{eq:proper-lc-in-p}.
\end{proof}

\begin{remark}
At \(\delta_i=0\) the leading coefficient again becomes a sum over the degree-maximising paths,
now confined to the block: at \(n=2\) with all \(r-1\) gaps zero it equals \(r/2\), the \(r\)
weakly increasing chains \(1\le\tau_1\le\dots\le\tau_{r-1}\le2\) each carrying a single factor
\(\tfrac12\). And the improper one-step window between saturation and threshold closes here: over
\(S_n^{(p)}\) the statistics satisfy \(\tp\le n\) and \(\indeg\ge2\) at every \(p\), so
\(\lvert\mathrm{FF}^{(p)}_{\mathrm{pr}}\rvert=\lvert S_n^{(p)}\rvert^{\,r}\) \textbf{exactly} when every
\(\delta_i\ge n-1\) (the \(p\)-analogue of Proposition \ref{prop-proper-threshold}), which is the
same point at which \(\kappa_i'\), hence the degree, saturates. The improper violating pair at
\(\delta_i=n\) was the full truncation against the void family, and the void family is not proper.

Verification, in the same style as for Theorem \ref{thm-asymptotics-p}: the degree and leading
coefficient by exact rational interpolation over \(125\) configurations with \(n\le5\), \(r\le4\)
and gaps \(0\) to \(n+2\) (degrees up to \(37\)), and the transfer formula
\eqref{eq:proper-p-product} against brute-force proper truncated sumsets for \(n\le3\), \(p\le5\),
\(r\le3\); no mismatches.
\label{rem-proper-p}
\end{remark}
\section{Verification, tables, and open problems}
\label{sec-verification}
\subsection{How the results were checked}
\label{sec:orgf77a9e1}

Every claim above labelled "verified" was checked against a route independent of the formula being
tested. Specifically:

\begin{itemize}
\item \textbf{Against brute force.} Direct enumeration of the order ideals of \(\mathcal P([n])\), forming the
sumset explicitly, with no recursion and no Kruskal--Katona: agreement in all \(41\)
configuration-and-variant pairs with \(n\le6\), \(r\le4\) and gaps \(0\) to \(6\), including
asymmetric gap vectors such as \((0,1,3)\) and \((0,2,4)\), for both the proper and the improper
count.
\item \textbf{Against Linusson's own tables.} All \(199\) printed entries of his Tables 5 and 6, exactly
(Remark \ref{rem-transcription}).
\item \textbf{Against a computer-algebra oracle.} Amata and Crupi's \texttt{ExteriorModules} package \autocite{AmataCrupi2021} for Macaulay2
enumerates lex submodules directly. Its counts \(43\), \(69\), \(99\) and \(236\) for \(n=3\) with
\(d=(0,0)\), \((0,1)\), \((-1,2)\) and \((0,0,1)\) are reproduced --- note that \((-1,2)\) exercises
the negative-degree hypothesis --- as is \(368\) for \(n=4\), \(d=(0,1)\).
\item \textbf{Against a proof assistant.} The combinatorial content of this paper has been formalised in Lean 4
with Mathlib, in the same development as the companion article's: the definitions and statements of
Sections \ref{sec-statistics} to \ref{sec-growth}, with the exceptions set out below. Every certified
statement rests on the three standard axioms of Lean's kernel (\texttt{propext}, \texttt{Classical.choice},
\texttt{Quot.sound}) and on no others, and the development contains no incomplete proof.

What follows is a description of what is \textbf{not} certified rather than a list of what is. That is the
more useful thing to publish: it is the boundary of the method, it says what the reader is being
asked to take on trust, and unlike an inventory of successes it does not go out of date as the
development grows. Three things lie outside.

\textbf{The algebra, in two quite different sizes.} None of Definition \ref{def-setup} is formalised — not
the exterior algebra, not monomial or lex submodules — and with it neither Lemma
\ref{lem-dictionary}, nor Lemma \ref{lem-lexproper}, nor the surjectivity half of Theorem
\ref{thm-chaining}. But these are not equally far away, and treating them alike would misdescribe
the work involved. Lemma \ref{lem-dictionary} is \textbf{elementary}: its two clauses are a spanning
observation and a dimension count, and what it lacks is not an argument but the objects — a
graded \(E\) with accessible graded pieces, monomial ideals, \(\mathfrak m^t\), and the
correspondence between monomial ideals and order ideals that makes \(\HF(E/I)\) an
\(f\)-vector. The surjectivity half of Theorem \ref{thm-chaining} needs all of that \textbf{and} Amata
and Crupi's theory of lex submodules — their Definition 2.6 \(\Leftrightarrow\) 2.7, the
existence and uniqueness of a lex representative in their Theorem 4.2, and Kruskal–Katona in
the exterior form of \textcite{AramovaHerzogHibi1997} — which is a substantially larger
undertaking, and is in any case open combinatorially. The formal development is
purely combinatorial: it counts order ideals and their \(f\)-vectors and never constructs a
submodule. Consequently, wherever a result needs the bridge from the combinatorics to the algebra,
that bridge is an explicit hypothesis rather than a proof. The counting identities of Theorems
\ref{thm-main} and \ref{thm-proper} are certified \textbf{unconditionally} as statements about chained tuples;
their readings as statements about \(\lvert\mathrm{FF}(\vec d)\rvert\) and
\(\lvert\mathrm{FF}_{\mathrm{pr}}(\vec d)\rvert\) are certified given that hypothesis, which
is exactly what "proved modulo" means in the text above.

\textbf{Linusson's results.} Theorem \ref{thm-linusson} and his Theorem 5.3 are cited rather than reproved,
so Lemma \ref{lem-joint} and everything resting on them is outside the formalisation: Theorem
\ref{thm-poly-p} and Theorems \ref{thm-asymptotics-p} and \ref{thm-asymptotics-p-proper}. For Lemma
\ref{lem-joint} this is deliberate and costs nothing, since \(A\) is defined by counting in
Definition \ref{def-A} and the recursion is nowhere needed.

\textbf{What "Perron root" is not.} Perron--Frobenius theory is not available in the library used. Of
Proposition \ref{prop-perron-threshold} what is certified is therefore the eigenvalue statement, for
both roots: above its threshold each of \(T(\delta)\) and \(T'(\delta)\) has rank one, \(\vec
  1\) is an eigenvector with eigenvalue \(\lvert\mathcal F_n\rvert\) and \(\lvert S_n\rvert\)
respectively, and no other eigenvalue is non-zero. The word \textbf{Perron} --- that this is the
eigenvalue of largest modulus, with a non-negative eigenvector --- is not certified, and in the
rank-one case the gap is only the comparison against \(0\), but it is a gap. The same limitation
leaves the growth claim of Proposition \ref{prop-recurrence} outside, and there, unlike here, no
degeneracy makes the gap small. That proposition's other half \textbf{is} certified: the linear recurrence
follows from the Cayley--Hamilton theorem, needs no spectral theory at all, and holds for both
matrices unconditionally.

Two kinds of thing were found by doing this. Several of the strengthenings recorded above were
found here, each of them a hypothesis that turned out to be unnecessary; and one hypothesis turned
out to be \textbf{missing} --- Theorem \ref{thm-proper}'s \(n\ge1\), without which the statement is false
at \(n=0\), where \(A'\) vanishes identically while
\(\lvert\mathrm{FF}_{\mathrm{pr}}\rvert=1\).

\item \textbf{Against previously recorded values.} The two sequences of Remark \ref{rem-proper-sequences}.
\item \textbf{Internal consistency.} The \(r=1\) specialisation (Corollary \ref{cor-r1}) and the product-regime
identities (Propositions \ref{prop-threshold} and \ref{prop-proper-threshold}) for \(n=2,\dots,6\),
\(r\le4\), both variants.
\end{itemize}

No mismatch was found in any of these.
\subsection{Open problems}
\label{sec:orgd8a4f6c}

\begin{enumerate}
\item \emph{Closed since the first draft.} The asymptotics are Theorem \ref{thm-asymptotics-p}: the degree
in \(p\) for every gap vector, and the leading coefficient whenever no gap vanishes, where the
summands decouple as \(F^p(n)\prod_iE^p(n,\kappa_i)\). The vanishing-gap case is
Remark \ref{rem-asymptotics-zero-gaps}. Both rest on Linusson's Theorem 5.3 and Proposition 5.4,
cited rather than reproved --- the same status Theorem \ref{thm-poly-p} already had. The proper
variant is Theorem \ref{thm-asymptotics-p-proper}: the same transport, on the principal block of
the proper transfer matrices, with \(E^p(n,1)\) in place of \(F^p(n)\) and
\(\kappa_i'=\max(1,n-\delta_i)\) in place of \(\kappa_i\).
\item \emph{Closed since the first draft.} The exact degrees at \(\delta=0\) are now
Theorem \ref{thm-delta0-degrees}, with explicit leading coefficients; the two structural entries of
\(A\) in Lemma \ref{lem-A-structure} are what the earlier draft was missing. The same argument
settled the two Perron-root observations, now Proposition \ref{prop-perron-threshold}. We leave the
item here, marked closed, rather than deleting it, since the first draft circulated.
\item \emph{Closed since the first draft.} The identity \(|S_n|=F^n(n-2)\) is now
Proposition \ref{prop-fnn2}: Alexander duality \(\Delta\mapsto\Delta^{\vee}\) is an involution of
\(\mathcal F_n\) exchanging the two statistics, under which properness becomes the vanishing of
the top two coordinates. Lemma \ref{lem-alexander} also explains the persymmetry of \(A\) and
yields the gap-reversal symmetry of Remark \ref{rem-gap-reversal}.
\item \textbf{A bijective proof.} Theorem \ref{thm-chaining} is proved by way of a canonical representative;
a direct bijection between \(\mathrm{FF}(\vec d)\) and the chained tuples, not routed through lex
submodules, would explain the constraint combinatorially rather than algebraically. Remark
\ref{rem-gap-reversal} is a constraint on any answer rather than a step toward one: such a
bijection must be compatible with the gap-reversing involution recorded there. The explicit
inverse of Remark \ref{rem-explicit-inverse} is half of an answer: it is closed-form, its
uniqueness is elementary, and it is equivariant under that involution. What is missing is a
direct proof that its clamps are \(f\)-vectors.
\item \textbf{Beyond constant gaps.} The formula handles an arbitrary sorted \(\vec d\), but the growth results of
Section \ref{sec-growth} assume a constant gap. Non-constant gap vectors give products of distinct
matrices. Proposition \ref{prop-normalform} is the structural answer, and it reduces the problem
to the pair \((B,P)\); the count also factorises across any saturated gap. What it does not give
is a closed form, and the natural routes to one are closed off: no two distinct \(T(\delta)\)
commute, they generate the full matrix algebra so there is no common triangularisation or
invariant filtration, and the only permutation symmetry is the reversal of
\eqref{eq:no-permutation-invariance}.
\end{enumerate}
\subsection{Integer sequences}
\label{sec:org4ba6c2b}

Six sequences arising here appear not to be in the OEIS: \(|S_n|\) itself, and the \(r=2\)
proper and improper counts at \(\delta=0\) and \(\delta=1\), together with one \(r=3\) example. Of
these \(|S_n|=1,2,5,16,70,457,4908,95248,3617645,286007155\), the number of \(f\)-vectors of
simplicial complexes on exactly \(n\) vertices --- seems the most natural candidate for submission,
being the first-difference sequence of A007695, which has no separate entry. Nothing has been
submitted.
\section*{Acknowledgements}
\label{sec:org7d5e35a}
Svante Linusson's 1999 paper is not merely cited here but used structurally: the refinement he
introduced for his own recursion turns out to be exactly the statistic the rank-\(r\) count requires,
and the present paper would have no content without it. Aliaksandra Kupreyeva's Bachelor's thesis \autocite{Kupreyeva2019}
(Linköping University, 2019; report LiTH-MAT-EX--2019/05--SE), written under my supervision, is where
counting questions about \(f\)-vectors entered this project.

Just as Geheimrat Goethe had the diligent Friedrich Wilhelm Riemer to help him write \emph{Die
Wahlverwandtschaften}, just as Grothendieck had Jean Dieudonné to help him complete EGA, I have
\textbf{Claude}! What exactly that amounted to is set out in the disclosure section below.

Computations were carried out in SageMath, in Python, and in Macaulay2, the latter using Amata and
Crupi's own \texttt{ExteriorModules} package.
\section*{Disclosure of AI assistance}
\label{sec:orgc0065da}
This manuscript was produced with substantial assistance from a large language model, Claude
(Anthropic), and the extent of it is set out here rather than left to be inferred. It is more
extensive in this paper than in its companion.

The model: identified that Linusson's refined count \(E^p(n,k)\) is the joint-distribution object the
rank-\(r\) problem needs, and proved Lemma \ref{lem-joint}; found and proved the chaining reformulation
of Theorem \ref{thm-chaining}, including the dictionary of Lemma \ref{lem-dictionary} and the hypotheses
recorded in Remark \ref{rem-sorted}; assembled and proved Theorem \ref{thm-main} and its consequences
(Corollary \ref{cor-r1}, Propositions \ref{prop-threshold} and \ref{prop-delta0}); found and proved
Lemma \ref{lem-lexproper}, without which the proper variant of Theorem \ref{thm-proper} would not follow;
derived the growth results of Section \ref{sec-growth}, including Theorem \ref{thm-poly-p}; transcribed
Linusson's Theorem 5.2 from the printed page and validated it against both of his appendix tables;
wrote all the supporting code; and drafted the exposition. The author set the problem, directed the
scope, and required the separation of proved from verified results that Section \ref{sec-intro}'s table
records.

Two corrections the model made to its own earlier work are worth recording, as evidence of what the
verification was for: an initial value of \(|S_9|\) was wrong by one, caught by checking against
OEIS A007695 and against Linusson's printed table; and a threshold in Propositions \ref{prop-threshold}
and \ref{prop-proper-threshold} had been stated for the wrong one of the two variants.

Every counting claim in this paper was checked along a route independent of the formula being tested,
as set out in Section \ref{sec-verification}, and the reference list has been checked against primary
sources. The author has verified the results to the best of his ability, made the final edits, and
assumes full responsibility for any errors, mistakes or gaps that remain.
\section*{Funding}
\label{sec:org34dea2f}
This research received no external funding. It was carried out as part of the author's regular duties
at Linköping University, within the fraction of that employment allocated to research.
\section*{Conflicts of interest}
\label{sec:org122a7d4}
The author declares no financial or non-financial conflicts of interest. The nature and extent of
large-language-model assistance in producing this manuscript is disclosed in full in the section
above.
\section*{References}
\label{sec:orgb50d301}
\printbibliography[heading=none]
\end{document}
